% Working copy based on paired version 2026-11-20-r2; bibliography maintained in ../refs.bib.
\documentclass[11pt,a4paper]{article}
\usepackage[T1]{fontenc}
\usepackage{lmodern}
\usepackage[margin=26mm,headheight=14pt]{geometry}
\usepackage{amsmath,amssymb,amsthm,mathtools}
\usepackage{microtype,booktabs,enumitem,xcolor}
\usepackage[nottoc,notlof,notlot,notindex]{tocbibind}
\usepackage[colorlinks=true,linkcolor=blue!45!black,citecolor=blue!45!black,urlcolor=blue!45!black]{hyperref}
\hypersetup{pdftitle={Uniform Fourth-Order Lam\'e and Stokes Regularity on Polyhedra with Applications to Strain-Gradient Elasticity},pdfsubject={Companion manuscript: polyhedral regularity, 20 November 2026}}
\newtheorem{theorem}{Theorem}[section]
\newtheorem{lemma}[theorem]{Lemma}
\newtheorem{proposition}[theorem]{Proposition}
\newtheorem{corollary}[theorem]{Corollary}
\theoremstyle{remark}
\newtheorem{remark}[theorem]{Remark}
\numberwithin{equation}{section}
\setlist[enumerate]{itemsep=3pt,topsep=5pt}
\setlength{\emergencystretch}{2em}
\allowdisplaybreaks[1]
\title{\bfseries Uniform Fourth-Order Lam\'e and Stokes Regularity on Polyhedra with Applications to Strain-Gradient Elasticity}
% Author names and affiliations to be supplied by the authors before submission.
\author{}
\date{20 November 2026}
\begin{document}
\maketitle
\begin{abstract}
We study clamped fourth-order Lam\'e and Stokes systems on bounded
connected polyhedral Lipschitz domains in three dimensions. On every
fixed convex polyhedron we establish an
\(H^{-1}\times L^2\)-to-\(H^3\times H^2\) shift for the fourth-order
Stokes system with compatible constraint data, together with
material-uniform fourth-order Lam\'e and classical elasticity estimates.
The vertex analysis retains the mean of the radial angular component
and absorbs its negative contribution through the divergence penalty.
A radial transformation of the test functions gives coercivity on the
full critical complex strip for every convex vertex cone, without an
additional solid-angle or spherical-eigenvalue restriction. A complete
mixed angular inverse, estimates at every axial edge frequency, and a
joint velocity--pressure Mellin shift yield the global Sobolev estimates.
Corrected reentrant-edge modes show that convexity is necessary, even
for \(L^4\) momentum loads in the fourth-order systems and the full
strain-gradient elasticity problem. As applications, the abstract
strain-gradient estimates of the companion planar paper give the three
weighted displacement scales \(\mu^{-1}\sqrt\iota\) and the three pressure
scales \([\lambda/(\lambda+2\mu)]\sqrt\iota\) for all
\(\lambda\ge0\), \(\mu>0\), and \(0<\iota<1\).
A polynomial load depending only on the polyhedron attains all six
scales simultaneously. The planar spectrum and the dimension-independent
divergence lifting and compatible-data results are cited from the
companion paper; the proofs here concern the three-dimensional
regularity and the polyhedral load construction.
\end{abstract}
\noindent\textbf{Keywords.} Strain gradient elasticity; convex polyhedra;
corner and edge singularities; fourth-order Stokes system; incompressible
limit; parameter-uniform regularity; boundary layers.

\smallskip
\noindent\textbf{2020 Mathematics Subject Classification.} Primary 35J58;
secondary 35J48, 35B65, 74B05.

\section{Introduction}

Clamped fourth-order elliptic systems on a polyhedron couple edge and
vertex singularities with the incompressible limit. We establish global
regularity for the fourth-order Lam\'e and Stokes systems on every fixed
convex polyhedron, and construct singular energy solutions at every
reentrant edge. Strain-gradient elasticity (SGE) provides an application
in which these regularity results must also control a vanishing internal
length and its associated boundary layer.

For constant coefficients, the operator studied here is
\[
 (I-\iota^2\Delta)\mathcal L_{\lambda,\mu},
 \qquad
 \mathcal L_{\lambda,\mu}u=-\mu\Delta u
                  -(\lambda+\mu)\nabla\operatorname{div}u,
\]
with \(u\in H^2_0(\Omega)^3\), \(\lambda\ge0\), \(\mu>0\), and
\(\iota>0\). We prove that convexity is both sufficient and necessary for
global \(H^3\) regularity for all \(L^2\) loads among bounded connected
polyhedral Lipschitz domains. Sufficiency includes every fixed convex
polyhedron. It imposes no additional bound on a dihedral angle, solid
angle, or first spherical eigenvalue.

The central issue is the fourth-order vertex system, for which scalar
regularity cannot simply be applied componentwise. Let \(\Sigma\) be the
spherical section of a pointed convex vertex cone. After rotation, its
closure lies in an open hemisphere. For the displacement exponent
\(s=\tfrac12+x+iy\), \(0\le x\le1\), and normalized penalty
\(\beta=(\lambda+\mu)/\mu\), we construct a radial test transform
\(T_{s,\beta}\) such that
\[
 \operatorname{Re}a_{\beta,s}(U,T_{s,\beta}U)
 \ge\frac{47}{192}\|\Delta_SU\|^2
       +\frac12y^4\|U\|^2.
\]
The estimate includes the entire closed complex strip and has a constant
independent of \(\beta\). Its proof uses a zero-mean hemisphere inequality
\(\|\Delta_Sv\|^2\ge10\|\nabla_Sv\|^2\), retains the radial mean in the
general case, and controls the resulting negative term by the
divergence penalty. In the Stokes limit the divergence constraint removes
that mean directly. This mechanism is what permits the full convex
polyhedral class.

Weak angular invertibility alone does not give the asserted Sobolev
shift. We construct the complete mixed inverse, including a
parameter-dependent right inverse of angular divergence and recovery of
the original pressure. We then prove estimates at every axial Fourier
frequency along an edge. At vertices a simultaneous Mellin shift moves
velocity and pressure from their energy line to the target line, and
dyadic annular estimates produce the unweighted global derivatives.
The argument keeps the correct negative force norms and the compatibility
of nonzero constraint data. A strong-space perturbation of the global
fourth-order Stokes inverse treats the nearly incompressible limit.

\paragraph{Relation to the planar companion paper.}
The companion paper~\cite{PolygonCompanion2026} proves the planar corner
spectrum and convexity criterion, the dimension-independent divergence
lifting and abstract SGE estimates, and the compatible-data convergence
criterion. We use those results as inputs, with their statements and
hypotheses recorded where needed. The new regularity arguments below
concern the spherical vertex pencil, the complete mixed angular inverse,
all axial Fourier frequencies at an edge, and the joint Mellin shift.
The polyhedral sharpness construction uses a three-dimensional
solenoidal field, and the necessity argument corrects an axial cutoff
to retain an integrable load. We retain these arguments in full.

\paragraph{Previous regularity hypotheses and uniform estimates.}
The need for a fourth-order clamped shift was explicit in the strain-gradient
finite element literature. Li, Ming and Wang \cite[Hypothesis~1]{LiMingWang2021Korn}
postulated an $H^{-1}$-to-$H^3$ estimate for the fourth-order Lam\'e
operator and observed that its validity on convex polytopes was unclear.
The convex-polyhedral theorem established here verifies this particular
hypothesis, and also supplies the dependence on the material parameters.
Chen, Huang and Huang \cite{ChenHuangHuang2023} developed parameter-robust estimates
under regularity assumptions and proved a two-dimensional fourth-order
Stokes shift for zero constraint data. The compatible-data planar extension and strong divergence lifting are
proved in~\cite[Propositions~2.8--2.9 and Lemma~2.7]{PolygonCompanion2026}
and are recalled as inputs below.

Liao, Ming and Xu \cite[Theorem~2.9]{LiaoMingXu2023} previously obtained
the half-order length-parameter upper bounds and the related higher-norm
estimates on smooth planar domains, with constants independent of
$\lambda$ and $\iota$. The abstract estimates and explicit material factors are established
in~\cite{PolygonCompanion2026}. Here we verify their regularity inputs
on convex three-dimensional polyhedra and prove matching
\emph{individual} lower bounds for six quantities using one domain-only
load, simultaneously for the full material and length parameter ranges.
The solenoidal square example of \cite{ChenHuangHuang2023} motivates this construction;
its additional gradient component makes the pressure scales nontrivial.
Liao and Ming \cite{LiaoMing2025} instead investigate a broken Hardy
inequality and an SGE formulation with natural boundary conditions; that
boundary-value problem is distinct from the present clamped one.

\paragraph{Implications for robust discretization.}
Huang, Huang and Tang \cite[Section~2.3, (2.7)--(2.8)]{HuangHuangTang2026robust}
use the same six weighted continuous estimates as regularity inputs for
a nonconforming SGE analysis on polytopal domains. The present results
verify these inputs on convex three-dimensional polyhedra, with a
sharper explicit pressure factor. Their separate classical-elasticity
$s$-regularity condition is stated for $2\le s\le3$; the present
classical estimate supplies its $s=2$ case, not all higher values of $s$.

Classical elasticity on convex nonsmooth domains has an extensive
theory. B\u{a}cu\c{t}\u{a} and Bramble \cite{BacutaBramble2003} prove
material-uniform \(H^2\) estimates on convex planar polygons.
Kozlov, Maz'ya and Schwab \cite{KozlovMazyaSchwab1992} analyze the
displacement pencil near a conical vertex; polyhedral elliptic and Stokes
regularity are developed further in \cite{MazyaRossmann2004,MazyaRossmann2007}.
We supply the classical estimate needed here with a proof that tracks
the full material range and the local constraint compatibility.
Our main contribution concerns the clamped fourth-order system and its
consequences for SGE, rather than the existence of classical fixed-parameter
convex-cone regularity.

Choi, Kim and Koh \cite{ChoiKimKoh2025} study the second-order Lam\'e
system on nonconvex planar polygons with one reentrant corner and
inhomogeneous Dirichlet data.
They obtain singular expansions, explicit singular coefficients, and
convergence of coefficients and regular remainders to the corresponding
Stokes quantities in the incompressible limit. The present pencil is
fourth order with clamped conditions, and its critical strip governs
three-dimensional \(H^3\) regularity. The results below include convex
polyhedral sufficiency, a converse for every nonconvex polyhedron,
and matching SGE estimates throughout the length and material parameter
ranges. Thus the two analyses concern related incompressible limits
but different operators, boundary conditions, and regularity conclusions.

The divergence liftings recalled from~\cite{PolygonCompanion2026} combine
the boundary-jet extension theory of Maz'ya, Mitrea and Shaposhnikova
\cite{MazyaMitreaShaposhnikova2010} with the support-preserving inverses of Costabel and
McIntosh~\cite{CostabelMcIntosh2010}. The planar clamped Lam\'e shift and
its corner-spectrum proof are also supplied by that companion paper,
using Costabel--Dauge's corner expansion theorem~\cite{CostabelDauge1994}.
Those proofs are not repeated here.
The baseline hemisphere inequalities used here are established in
Mayboroda and Maz'ya \cite[Lemma~2.3, (2.16), (2.25)]{MayborodaMazya2008};
the strengthened zero-mean inequality needed for the coupled radial
test is proved below.

The main statements and their quantifiers appear in Section~\ref{sec:main}.
Sections~\ref{sec:hemispherical-pencil}--\ref{sec:g-lameuniform}
establish fourth-order regularity. Section~\ref{sec:est-classical}
proves the classical material-uniform input.
Sections~\ref{sec:est-sge} and \ref{sec:est-sharpness}
give the SGE upper and lower bounds.
Section~\ref{sec:necessity} proves the converse by a localized edge
construction. Section~\ref{comp:section} records a compatible-data
consequence of the companion paper and its explicit material scaling.

\section{Formulation and main results}\label{sec:main}

\subsection{Spaces, operators, and energy solutions}

A polyhedral Lipschitz domain means a bounded connected open subset of
\(\mathbb R^3\) with Lipschitz boundary consisting of finitely many
planar faces. Coplanar subdivisions do not create geometric edges.
All material coefficients are spatially constant. The space
\(H^k_0(\Omega)\) is the closure of \(C_c^\infty(\Omega)\) in \(H^k\);
in particular \(H^2_0\) imposes zero displacement and zero first-order
traces. We use \(H^{-k}(\Omega)=(H^k_0(\Omega))^*\),
\(L^2_0(\Omega)=\{q\in L^2:\int_\Omega q=0\}\), and the usual full
norms \(\|\cdot\|_{H^k}\) and seminorms \(|\cdot|_{H^k}\).
An omitted domain is \(\Omega\). Vector and tensor norms are componentwise.
The real spaces are complexified in the transform arguments, with
inner products linear in the first variable.

Write
\[
 \epsilon(u)=\tfrac12(\nabla u+\nabla u^\top),\qquad
 \sigma_{\lambda,\mu}(u)=2\mu\epsilon(u)
                      +\lambda(\operatorname{div}u)I,
\]
\[
 P_{\lambda,\mu}=-\Delta\mathcal L_{\lambda,\mu}
 =\mu\Delta^2+(\lambda+\mu)\nabla\Delta\operatorname{div}.
\]
The symmetric forms are
\begin{align}
 A_{\lambda,\mu}(u,v)
 &=2\mu(\epsilon(u),\epsilon(v))
                  +\lambda(\operatorname{div}u,\operatorname{div}v),\\
 C_{\lambda,\mu}(u,v)
 &=2\mu(\nabla\epsilon(u),\nabla\epsilon(v))
                  +\lambda(\nabla\operatorname{div}u,\nabla\operatorname{div}v).
\end{align}
Integration by parts and density give
\[
 A_{\lambda,\mu}(u,u)
 =\mu\|\nabla u\|^2+(\lambda+\mu)\|\operatorname{div}u\|^2
 \quad(u\in H^1_0),
\]
\[
 C_{\lambda,\mu}(u,u)
 =\mu\|D^2u\|^2+(\lambda+\mu)\|\nabla\operatorname{div}u\|^2
 \quad(u\in H^2_0).
\]
Consequently the clamped fourth-order problem is coercive on \(H^2_0\).
For \(f\in L^2(\Omega)^3\), the unique SGE and classical energy solutions
are defined by
\begin{align}
 A_{\lambda,\mu}(u_\iota,v)+\iota^2C_{\lambda,\mu}(u_\iota,v)
  &=(f,v) &&(v\in H^2_0(\Omega)^3),\label{eq:main-sge}\\
 A_{\lambda,\mu}(u_0,v)&=(f,v)
               &&(v\in H^1_0(\Omega)^3).\label{eq:main-classical}
\end{align}
These are the weak forms of
\((I-\iota^2\Delta)\mathcal L_{\lambda,\mu}u_\iota=f\) and
\(\mathcal L_{\lambda,\mu}u_0=f\).
We use the pressure convention
\[
 p_\iota=\lambda\operatorname{div}u_\iota,\qquad
 p_0=\lambda\operatorname{div}u_0,\qquad
 \alpha=\frac{\lambda}{\lambda+2\mu}.
\]
The auxiliary mixed pressures used in the regularity proofs are
normalized differently and are always specified where introduced.

For the fourth-order Stokes system let
\[
 -\Delta\chi_\Omega=1,\quad \chi_\Omega\in H^1_0(\Omega),\qquad
 Q(\Omega)=H^1_0(\Omega)\cap L^2_0(\Omega),\qquad
 L^2_\chi(\Omega)=\{g\in L^2:(g,\chi_\Omega)=0\}.
\]
Its equations are
\begin{equation}\label{eq:main-stokes}
 \Delta^2v+\nabla\Delta q=F,\qquad
 \Delta\operatorname{div}v=g,\qquad
 (v,q)\in H^2_0(\Omega)^3\times Q(\Omega).
\end{equation}
The condition \((g,\chi_\Omega)=0\) follows from
\(\int_\Omega\operatorname{div}v=0\); it is also sufficient for energy
solvability. The pressure has zero Dirichlet trace and zero mean,
but no normal derivative boundary condition.

\subsection{Global shifts and the geometric criterion}

\begin{theorem}[Material-uniform regularity]\label{thm:main-regularity}
Let \(\Omega\) be a fixed bounded convex polyhedron.
\begin{enumerate}[label=\textup{(\roman*)}]
\item For \(F\in H^{-1}(\Omega)^3\) and \(g\in L^2_\chi(\Omega)\),
the solution of \eqref{eq:main-stokes} satisfies
\begin{equation}\label{eq:main-stokes-shift}
 \|v\|_{H^3}+\|q\|_{H^2}
 \le C_\Omega(\|F\|_{H^{-1}}+\|g\|_{L^2}).
\end{equation}
\item If \(P_{\lambda,\mu}u=F\in H^{-1}(\Omega)^3\) with \(u\in H^2_0\),
then
\begin{equation}\label{eq:main-lame4}
 \mu\|u\|_{H^3}+(\lambda+\mu)\|\operatorname{div}u\|_{H^2}
 \le C_\Omega\|F\|_{H^{-1}}.
\end{equation}
\item The solution of \eqref{eq:main-classical} satisfies
\begin{equation}\label{eq:main-lame2}
 \mu\|u_0\|_{H^2}+(\lambda+\mu)\|\operatorname{div}u_0\|_{H^1}
 \le C_\Omega\|f\|_{L^2}.
\end{equation}
\end{enumerate}
The constants in \eqref{eq:main-lame4}--\eqref{eq:main-lame2}
are independent of all \(\lambda\ge0\) and \(\mu>0\).
\end{theorem}

The three parts are proved in Theorems~\ref{thm:g-stokes},
\ref{thm:g-lameuniform}, and \ref{thm:est-classical}, respectively.
Constants throughout may depend on the fixed domain. No uniformity over
degenerating families of polyhedra is asserted.

\begin{theorem}[Convexity criterion]\label{thm:main-convexity}
Let \(\Omega\) be a bounded connected polyhedral Lipschitz domain.
For each fixed \(\lambda\ge0\), \(\mu>0\), and \(\iota>0\), the following
properties are equivalent:
\begin{enumerate}[label=\textup{(\roman*)}]
\item \(\Omega\) is convex.
\item For every \(F\in H^{-1}(\Omega)^3\), the clamped energy solution
of \(P_{\lambda,\mu}u=F\) belongs to \(H^3(\Omega)^3\).
\item For every \(f\in L^2(\Omega)^3\), the SGE solution
\(u_\iota\) belongs to \(H^3(\Omega)^3\).
\item For every \(F\in H^{-1}(\Omega)^3\), the solution of
\eqref{eq:main-stokes} with \(g=0\) belongs to
\(H^3(\Omega)^3\times H^2(\Omega)\).
\item For every \(f\in L^2(\Omega)^3\), the classical solution
\(u_0\) belongs to \(H^2(\Omega)^3\).
\end{enumerate}
On a nonconvex domain, the counterexamples in \textup{(ii)--(iv)}
can be chosen with \(L^4\) momentum loads. The counterexample loads may
depend on the fixed material and length parameters.
\end{theorem}

The sufficiency is contained in Theorem~\ref{thm:main-regularity}.
Section~\ref{sec:necessity} proves the converse and explains the reduction
to a reentrant edge. Convexity is a condition for regularity for
\emph{every} load; a nonconvex domain can of course admit smooth solutions
for particular loads.

\subsection{SGE applications: six sharp scales}

For the solutions associated with one load \(f\), define
\[
 \begin{array}{lll}
 T_1=|u_\iota-u_0|_{H^1},&
 T_2=\iota\|u_\iota\|_{H^2},&
 T_3=\iota^2\|u_\iota\|_{H^3},\\
 T_4=\|p_\iota-p_0\|_{L^2},&
 T_5=\iota|p_\iota|_{H^1},&
 T_6=\iota^2\|p_\iota\|_{H^2}.
 \end{array}
\]
\begin{theorem}[Uniform bounds and simultaneous optimality]
\label{thm:main-sharp}
For every fixed bounded convex polyhedron and every \(f\in L^2(\Omega)^3\),
\begin{align}
 T_1+T_2+T_3&\le C_\Omega\mu^{-1}\sqrt{\iota}\,\|f\|_{L^2},
 \label{eq:main-u}\\
 T_4+T_5+T_6&\le C_\Omega\alpha\sqrt{\iota}\,\|f\|_{L^2},
 \label{eq:main-p}
\end{align}
for all \(\lambda\ge0\), \(\mu>0\), and \(0<\iota<1\).
There is one nonzero polynomial \(f_\Omega\), depending only on \(\Omega\),
such that for this same load
\[
 c_\Omega\mu^{-1}\sqrt{\iota}\le T_j
       \le C_\Omega\mu^{-1}\sqrt{\iota}\quad(j=1,2,3),
\]
\[
 c_\Omega\alpha\sqrt{\iota}\le T_j
       \le C_\Omega\alpha\sqrt{\iota}\quad(j=4,5,6)
\]
throughout that entire parameter range.
When \(\lambda=0\), the three pressure quantities vanish identically.
\end{theorem}

The upper bounds follow from Proposition~\ref{prop:est-sge-transfer},
whose regularity inputs have all been proved on this domain class.
The construction and the six separate lower bounds are given in
Theorem~\ref{thm:est-sharpness}.
These are statements about the continuous length parameter \(\iota\),
not about a discretization parameter.

The compatible-data criterion is recalled, with attribution, in
Section~\ref{comp:section}. It is a dimension-independent input from
\cite{PolygonCompanion2026}, rather than an additional polyhedral
regularity theorem.

\section{Uniform weak pencils on hemispherical cone sections}
\label{sec:hemispherical-pencil}

All inner products in this section are linear in their first argument.
Let $\mathcal H=\{\omega\in\mathbb S^2:\omega_3>0\}$ and
$M=|\mathcal H|=2\pi$. We write $z=\omega_3$ and, for a scalar function $v$,
\[
 N_v=\|v\|_{L^2(\mathcal H)}^2,\qquad
 A_v=\|\nabla_Sv\|_{L^2(\mathcal H)}^2,\qquad
 D_v=\|\Delta_Sv\|_{L^2(\mathcal H)}^2,\qquad
 J_v=M^{-1}\left|\int_{\mathcal H}v\right|^2.
\]
The subscript is omitted when no ambiguity can arise. Vector derivatives
are applied to Cartesian components. All the inequalities below hold for
complex-valued functions.

\subsection{Two hemisphere inequalities and a mean-zero improvement}

\begin{lemma}[Hemisphere inequalities]\label{lem:hemisphere-basic20}
For $v\in H_0^1(\mathcal H)$,
\begin{equation}\label{eq:hemisphere-poincare20}
 A_v\ge2N_v,\qquad A_v\ge6(N_v-J_v).
\end{equation}
For $v\in H_0^2(\mathcal H)$,
\begin{equation}\label{eq:hemisphere-buckling20}
 D_v\ge6A_v.
\end{equation}
If, in addition, $\int_{\mathcal H}v=0$, then
\begin{equation}\label{eq:hemisphere-zero-mean20}
 A_v\ge6N_v,\qquad D_v\ge10A_v.
\end{equation}
\end{lemma}

\begin{proof}
The first Dirichlet eigenvalue on the hemisphere is $2$, with eigenfunction
$z$, proving the first inequality in \eqref{eq:hemisphere-poincare20}.
The remaining baseline inequalities in \eqref{eq:hemisphere-poincare20}
and \eqref{eq:hemisphere-buckling20} are given by
Mayboroda--Maz\'ya \cite[Lemma~2.3, especially (2.16) and (2.25)]{MayborodaMazya2008}.
For clarity about the exact clamped and mean conditions, we reproduce
their short spectral argument before proving the additional zero-mean
$D\ge10A$ estimate.
Decompose in azimuthal modes $e^{im\varphi}$, $m\in\mathbb Z$. The quantities
$N,A,D$ split orthogonally. In the axisymmetric mode, even reflection across
the equator is an $H^1$ function on the sphere. Its spherical harmonic
expansion contains the even Legendre polynomials $P_0,P_2,\ldots$; the first
nonconstant eigenvalue is $6$. For $|m|\ge1$, odd reflection of the zero
Dirichlet trace gives degrees $\ell$ satisfying $\ell+|m|$ odd. Hence
$\ell\ge|m|+1$ and the gradient quotient is at least
$(|m|+1)(|m|+2)\ge6$. The constant component has squared norm $J$.
This proves $A\ge6(N-J)$.

If $v\in H_0^2(\mathcal H)$, then $\int_{\mathcal H}\Delta_Sv=0$.
With $v_{\mathcal H}=M^{-1}\int_{\mathcal H}v$, integration by parts yields
\[
 A=-(\Delta_Sv,v-v_{\mathcal H}),\qquad
 A^2\le D(N-J)\le DA/6.
\]
This proves \eqref{eq:hemisphere-buckling20}. The first mean-zero estimate
is already contained in \eqref{eq:hemisphere-poincare20}.

It remains to show $D\ge10A$ for mean-zero $v\in H_0^2(\mathcal H)$.
The zero normal derivative permits even $H^2$ reflection. Put
$\lambda_\ell=\ell(\ell+1)$.

For the axisymmetric mode, the mean-zero condition removes degree zero,
and, suppressing the common factor $2\pi$, we have
\[
 v=\sum_{\substack{\ell\ge2\\\ell\text{ even}}}a_\ell P_\ell(z),
 \qquad \int_0^1P_\ell(z)^2\,dz=(2\ell+1)^{-1}.
\]
Consequently
\[
 D-12A=-\frac{36}{5}|a_2|^2+E_0,
 \qquad
 E_0=\sum_{\substack{\ell\ge4\\\ell\text{ even}}}
 \frac{\lambda_\ell(\lambda_\ell-12)}{2\ell+1}|a_\ell|^2.
\]
The equator trace $v(0)=0$ and $P_2(0)=-1/2$ imply
\[
 \frac{|a_2|^2}{4}\le E_0S_0,
 \qquad
 S_0=\sum_{\substack{\ell\ge4\\\ell\text{ even}}}
 \frac{(2\ell+1)|P_\ell(0)|^2}{\lambda_\ell(\lambda_\ell-12)}.
\]
The identity $P_{2n}(0)=(-1)^n\binom{2n}{n}4^{-n}$ gives
$|P_\ell(0)|^2\le9/64$ for even $\ell\ge4$, while
$\lambda_\ell-12\ge(2/5)\lambda_\ell$. Therefore
\[
 S_0\le\frac{45}{128}\sum_{\ell=4}^\infty
 \left(\frac1{\ell^2}-\frac1{(\ell+1)^2}\right)
 =\frac{45}{2048}<\frac5{144}.
\]
It follows that $E_0\ge(36/5)|a_2|^2$, so $D\ge12A$ on this mode.

For each mode $|m|=1$, even reflection permits odd degrees. Use
\[
 P_\ell^1(z)=-(1-z^2)^{1/2}P_\ell'(z),\qquad
 \int_0^1|P_\ell^1|^2\,dz=\frac{\lambda_\ell}{2\ell+1},\qquad
 P_\ell^1(0)=-\ell P_{\ell-1}(0).
\]
Writing $v=\sum_{\ell\text{ odd}}a_\ell P_\ell^1$ gives
\[
 D-10A=-\frac{32}{3}|a_1|^2+E_1,
 \qquad
 E_1=\sum_{\substack{\ell\ge3\\\ell\text{ odd}}}
 \frac{\lambda_\ell^2(\lambda_\ell-10)}{2\ell+1}|a_\ell|^2.
\]
Since $P_1^1(0)=-1$, the zero equator trace yields
\[
 |a_1|^2\le E_1S_1,
 \qquad
 S_1=\sum_{\substack{\ell\ge3\\\ell\text{ odd}}}
 \frac{(2\ell+1)|P_{\ell-1}(0)|^2}
      {(\ell+1)^2(\lambda_\ell-10)}.
\]
The term $\ell=3$ equals $7/128$. For odd $\ell\ge5$,
$|P_{\ell-1}(0)|^2\le9/64$ and
$\lambda_\ell-10\ge(2/3)\lambda_\ell$. Moreover,
\[
 \frac{2\ell+1}{\ell(\ell+1)^3}
 \le\frac1{\ell^2}-\frac1{(\ell+1)^2},
\]
because the ratio of the left side to the right side is $\ell/(\ell+1)$.
Thus
\[
 S_1\le\frac7{128}+\frac{27}{128}
 \sum_{\ell=5}^\infty
 \left(\frac1{\ell^2}-\frac1{(\ell+1)^2}\right)
 =\frac{101}{1600}<\frac3{32}.
\]
This gives $E_1\ge(32/3)|a_1|^2$ and proves $D\ge10A$ on these modes.

Finally, for $|m|\ge2$, odd $H^1$ reflection gives $A\ge12N$.
The inequality $A^2\le DN$ then gives $D\ge12A$.
Summing the orthogonal components proves \eqref{eq:hemisphere-zero-mean20}.
The expansions converge in $H^2(\mathbb S^2)$ after even reflection.
The convergent series $S_0,S_1$ bound the equator trace
expansions by weighted Cauchy--Schwarz; their finite partial sums
are uniformly bounded, so the limiting traces are well defined. Alternatively,
one first argues for smooth clamped functions and then uses $H_0^2$
density, correcting the mean with a fixed interior smooth function.
\end{proof}

\subsection{The weak angular operator and a radial test transform}

Let $\Sigma$ be a Lipschitz angular subdomain whose closure is contained in
an open hemisphere, rotated to $\mathcal H$. Zero extension maps
$H_0^k(\Sigma)$ to $H_0^k(\mathcal H)$ for $k=1,2$. The hemisphere
inequalities therefore apply to functions on $\Sigma$. A pointed convex
polyhedral cone has such a section.

Set
\[
 V=H_0^2(\Sigma;\mathbb C^3),\qquad
 s=\frac12+x+iy,\quad 0\le x\le1,\qquad
 t=1-\overline s.
\]
For $U\in V$, write
\[
 R=U\cdot\omega,\qquad W=U-R\omega,\qquad
 d_sU=\operatorname{div}_S W+(s+2)R=:h.
\]
Here $W$ is tangential and $h\in H_0^1(\Sigma)$. Define
\[
 L_\zeta=\Delta_S+\zeta(\zeta+1),\qquad
 C_s(h,k)=(\nabla_Sh,\nabla_Sk)-s(s-1)(h,k).
\]
For smooth angular fields, the homogeneous ansatz gives
\begin{equation}\label{eq:strong-angular-pencil23}
 \begin{split}
 &(\Delta^2+\beta\nabla\Delta\operatorname{div})(r^sU)
       =r^{s-4}\mathfrak A_\beta(s)U,\\
 &\mathfrak A_\beta(s)U
       =L_{s-2}L_sU
        +\beta\bigl(\nabla_S+(s-3)\omega\bigr)L_{s-1}d_sU.
 \end{split}
\end{equation}
Indeed $\Delta(r^sU)=r^{s-2}L_sU$ and
$\operatorname{div}(r^sU)=r^{s-1}d_sU$, while the gradient of
$r^{s-3}k$ equals $r^{s-4}(\nabla_Sk+(s-3)k\omega)$.
Angular integration by parts yields the normalized weak pencil
\begin{equation}\label{eq:weak-pencil20}
 a_{\beta,s}(U,V)=(L_sU,L_tV)+\beta C_s(d_sU,d_tV),\qquad\beta>0.
\end{equation}
The conjugate coefficient of $L_t$ is $(s-1)(s-2)$, and
$\overline t+2=3-s$, so the first-order angular factor in
\eqref{eq:strong-angular-pencil23} pairs with $-d_tV$.
Consequently
\[
 \langle\mathfrak A_\beta(s)U,V\rangle
 =(L_sU,L_tV)-\beta\langle L_{s-1}d_sU,d_tV\rangle
 =a_{\beta,s}(U,V)
\]
on $V\times V$ by density; the middle pairing is between $H^{-1}$ and
$H^1_0$.
Thus \eqref{eq:strong-angular-pencil23} also defines a distributional
map $V\to V^*$; no additional angular regularity is assumed.
In particular the form depends analytically on $s$ when represented on the
fixed spaces $V$ and its antidual.

For a scalar function $R$ put $a=s(s+1)$ and
\[
 B_R^0=D_R-2(a+1)A_R+a(a-2)N_R,
 \qquad B_0=(L_sU,L_tU).
\]
The key identity is
\begin{equation}\label{eq:radial-cross20}
 (L_sU,L_t(R\omega))=B_R^0+4C_s(h,R).
\end{equation}
Here and below scalar norms can be computed on $\Sigma$ or after zero
extension to $\mathcal H$.

To verify \eqref{eq:radial-cross20}, let $\Delta_{\mathrm{cov}}$ denote the
connection Laplacian on tangential fields, with the same sign convention
as $\Delta_S$. The radial and tangential parts of $L_sU$ are
\[
 \Delta_SR+(a+2s+2)R-2h,
 \qquad \Delta_{\mathrm{cov}}W+(a-1)W+2\nabla_SR.
\]
Using the Ricci curvature $1$ of $\mathbb S^2$ gives
\[
 (\Delta_{\mathrm{cov}}W,\nabla_SR)
 =-(\operatorname{div}_S W,\Delta_SR+R).
\]
Substitute $\operatorname{div}_S W=h-(s+2)R$ and integrate by parts.
The coefficients of $A_R,N_R$ become $-2(a+1),a(a-2)$, and the
remaining terms are $4(\nabla_Sh,\nabla_SR)-4s(s-1)(h,R)$.
This proves the identity for smooth clamped functions and then for $V$
by density.

Let $\delta=4/\beta$ and introduce the bounded local multiplier
\begin{equation}\label{eq:test-transform20}
 T_{s,\beta}U=W+\frac{s+2+\delta}{3-\overline s+\delta}R\omega.
\end{equation}
Direct calculation gives
\[
 d_tT_{s,\beta}U=h-\frac{2x\delta}{3-\overline s+\delta}R.
\]
Since the second slot in \eqref{eq:weak-pencil20} is antilinear,
\eqref{eq:radial-cross20} and $\beta\delta=4$ cancel the mixed terms
exactly:
\begin{equation}\label{eq:exact-transformed20}
 a_{\beta,s}(U,T_{s,\beta}U)
 =B_0+\frac{2x}{3-s+\delta}B_R^0+\beta C_s(h,h).
\end{equation}

\subsection{The radial block with its mean retained}

Put $Y=y^2$, $\alpha=x^2+2x+3/4$, and
\[
 C_\alpha(R)=D_R-2(\alpha+1)A_R+\alpha(\alpha-2)N_R,
 \qquad J=M^{-1}\left|\int_{\mathcal H}R\right|^2.
\]
For $0\le x\le1/2$, Lemma~\ref{lem:hemisphere-basic20} implies
\begin{equation}\label{eq:radial-low20}
 C_\alpha(R)\ge\frac12(\alpha-2)(\alpha-4)A_R\ge0.
\end{equation}
For $1/2<x\le1$, set $\eta=z^2$ and
\[
 c=3M^{-1}\int_{\mathcal H}R,\qquad R_0=R-c\eta.
\]
Then $R_0\in H_0^2(\mathcal H)$, $\int_{\mathcal H}R_0=0$, and
$|c|^2M=9J$. It is immaterial that $R_0$ need not be supported in
$\Sigma$: this decomposition is used only for a scalar comparison estimate.
Since $2<\alpha\le15/4$, \eqref{eq:hemisphere-zero-mean20} yields
\begin{equation}\label{eq:radial-zero-mean20}
 C_\alpha(R_0)\ge(\alpha^2-14\alpha+48)N_{R_0}
 \ge\frac{153}{16}N_{R_0}.
\end{equation}
Elementary integration gives
\[
 \Delta_S\eta=2-6\eta,\quad
 N_\eta=M/5,\quad A_\eta=8M/15,\quad D_\eta=16M/5,
\]
\[
 \|\eta-1/3\|^2=4M/45,\qquad
 M^{-1}C_\alpha(\eta)=\frac{3\alpha^2-22\alpha+32}{15}\ge-\frac59.
\]
The polarized cross term is
\[
 C_\alpha(R_0,\eta)=(\alpha-2)(\alpha-12)(R_0,\eta),
\]
whose coefficient has absolute value at most $231/16$.
Completing the square in $\sqrt{N_{R_0}}$ proves
\begin{equation}\label{eq:radial-mean-loss20}
 C_\alpha(R)\ge
 -|c|^2M\left(\frac59+
 \frac{(231/16)^2(4/45)}{153/16}\right)
 =-\frac{7629}{340}J\ge-\frac{45}{2}J.
\end{equation}

Write $d_\delta=5/2-x+\delta$. Expansion of the complex radial term gives
\begin{align}
 |3-s+\delta|^2\operatorname{Re}\frac{B_R^0}{3-s+\delta}
 ={}&d_\delta C_\alpha(R)
 +Y\bigl((9+2x+2\delta)A_R-[M_0(x)+\delta Q(x)]N_R\bigr)
 \notag\\
 &+(13/2+3x+\delta)Y^2N_R,\label{eq:radial-complex20}\\
 M_0(x)={}&-2x^3+15x^2+(67/2)x+31/4,
 \qquad Q(x)=6x^2+12x+7/2.\notag
\end{align}
The last term is nonnegative. The mean-sensitive bound
$A_R\ge6(N_R-J)$ shows that the expression multiplying $Y$ is at least
\[
 (47/4-19\delta/2)N_R-6(9+2x+2\delta)J.
\]
Indeed $6(9+2x)-M_0(x)$ decreases on $[0,1]$ to $47/4$, and
$12-Q(x)\ge-19/2$. After multiplication by
$2xY/(d_\delta^2+Y)$, the first term loses at most $2YN_R$ because
\[
 2(\delta+3/2)^2-(19\delta-47/2)
 =2\delta^2-13\delta+28>0.
\]
The mean loss in the $Y$ term is
\[
 L_YJ,\qquad L_Y=\frac{12x(9+2x+2\delta)Y}{d_\delta^2+Y}.
\]
The constant radial term loses no more than $L_CJ$, where
$L_C=0$ for $x\le1/2$ and $L_C\le45/(\delta+3/2)$ otherwise.
With $B=\operatorname{Re}B_0$ and $E_h=\operatorname{Re}C_s(h,h)$,
\eqref{eq:exact-transformed20} therefore yields
\begin{equation}\label{eq:before-mean-absorption20}
 \operatorname{Re}a_{\beta,s}(U,T_{s,\beta}U)
 \ge B-2YN_R+\beta E_h-(L_C+L_Y)J.
\end{equation}

\begin{lemma}[Absorption of the radial mean]\label{lem:radial-mean-absorption22}
For every $U\in V$, $\beta>0$ and $s=\tfrac12+x+iy$ with $0\le x\le1$,
let $R=U\cdot\omega$ and let $T_{s,\beta}$ be given by
\eqref{eq:test-transform20}. Then
\begin{equation}\label{eq:after-mean-absorption20}
\begin{aligned}
 \operatorname{Re}a_{\beta,s}(U,T_{s,\beta}U)
 &\ge \operatorname{Re}(L_sU,L_tU)-2y^2\|R\|^2\\
 &\ge \operatorname{Re}(L_sU,L_tU)-2y^2\|U\|^2.
\end{aligned}
\end{equation}
\end{lemma}
\begin{proof}
By \eqref{eq:before-mean-absorption20}, it suffices to show
$\beta E_h\ge(L_C+L_Y)J$.
Since $W$ has zero trace,
\[
 \int_{\mathcal H}h=(s+2)\int_{\mathcal H}R,
 \qquad N_h\ge|s+2|^2J,
 \qquad E_h\ge(9/4-x^2+Y)N_h.
\]
If $1/2<x\le1$, these inequalities give
\[
 \beta E_h\ge\delta^{-1}(45+41Y+4Y^2)J.
\]
The constant part absorbs $L_CJ$, and the $Y$ part absorbs $L_YJ$ because
\[
 L_Y\le\frac{12(11+2\delta)Y}{(\delta+3/2)^2+Y}
 \le\frac{41Y}{\delta},
\]
where the second inequality follows from
$17\delta^2-9\delta+369/4>0$.
If $0\le x\le1/2$, then $L_C=0$ and
\[
 L_Y\le\frac{6(10+2\delta)Y}{(\delta+2)^2+Y}
 \le\frac{33Y}{\delta},\qquad
 \beta E_h\ge\delta^{-1}(50+33Y+4Y^2)J.
\]
Here the middle inequality follows from
$21\delta^2+72\delta+132>0$. These bounds prove the claimed estimate
for every $\delta>0$.
\end{proof}

The preceding calculation separates the geometric issue from the material
parameters. The possible negative part of the radial block is bounded
only by $J$, the squared hemispherical mean in
\eqref{eq:radial-mean-loss20}, rather than by the full $L^2$ norm of
$R$. The exact divergence identity
$\int_{\mathcal H}h=(s+2)\int_{\mathcal H}R$ then absorbs this loss
with the penalty term, giving \eqref{eq:after-mean-absorption20}.
Thus the hemispherical comparison applies to every convex vertex
cone without an additional vertex-dependent spectral condition.
The numerical coercivity constants in this argument are not asserted
to be optimal.

\subsection{Coercivity and the complete weak inverse}

\begin{theorem}[Uniform weak angular inverse]\label{thm:weak-angular20}
For every $\beta>0$, every $s$ with $1/2\le\operatorname{Re}s\le3/2$,
and every $U\in V$, the transform \eqref{eq:test-transform20} satisfies
\begin{equation}\label{eq:weak-angular-coercivity20}
 \operatorname{Re}a_{\beta,s}(U,T_{s,\beta}U)
 \ge\frac{47}{192}\|\Delta_SU\|^2+\frac12y^4\|U\|^2.
\end{equation}
For $\kappa=(1+y^2)^{1/2}$ equip $V$ with the parameter norm
\[
 \|U\|_{V_y}^2=\|\Delta_SU\|^2+
 \kappa^2\|\nabla_SU\|^2+\kappa^4\|U\|^2.
\]
The weak pencil is an isomorphism $V_y\to V_y^*$, and
\begin{equation}\label{eq:weak-angular-inverse20}
 \|U\|_{V_y}\le C\|a_{\beta,s}(U,\cdot)\|_{V_y^*},
\end{equation}
with $C$ independent of $s$ and $\beta$.
\end{theorem}

\begin{proof}
Put $D=\|\Delta_SU\|^2$, $A=\|\nabla_SU\|^2$, $N=\|U\|^2$ and
\[
 \underline B=D-\frac72A-\frac{15}{16}N
   +Y\left(2A-\frac72N\right)+Y^2N.
\]
The scalar expansion of $B$ gives
\[
 B-\underline B=(1-x^2)\{2A+(3/2-x^2+6Y)N\}\ge0.
\]
Applying Lemma~\ref{lem:hemisphere-basic20} componentwise gives
$D\ge6A$ and $A\ge2N$. The exact decomposition
\begin{align*}
 \underline B-2YN={}&\frac{47}{192}D+\frac12Y^2N
 +\frac{145}{192}(D-6A)
 +\left(\frac{33}{32}+2Y\right)(A-2N)
 +\frac12(Y-3/2)^2N
\end{align*}
and \eqref{eq:after-mean-absorption20} prove
\eqref{eq:weak-angular-coercivity20}.

The radial multiplier in $T_{s,\beta}$ has modulus between $1$ and
$7/3$. Its inverse is bounded, and multiplication by the smooth projection
$\omega\otimes\omega$ preserves $H_0^2$. The product rule therefore bounds
$T_{s,\beta}$ and its inverse uniformly in the parameter norms.
Since $D\ge6A$ and $A\ge2N$, integration by parts and Cauchy--Schwarz give
\[
 A\le D/6,\qquad N\le D/12,\qquad
 YA\le\sqrt{DY^2N}\le\tfrac12(D+Y^2N).
\]
Using also $2YN\le YA$, we obtain
\[
 \|U\|_{V_y}^2=D+(1+Y)A+(1+Y)^2N
 \le\tfrac94D+2Y^2N.
\]
Thus the right side of~\eqref{eq:weak-angular-coercivity20} is at least
$(47/432)\|U\|_{V_y}^2$.

For each fixed finite $\beta$, the form
$(U,V)\mapsto a_{\beta,s}(U,T_{s,\beta}V)$ is continuous and coercive on
$V_y$. Complex Lax--Milgram gives a unique solution for every functional
$F\circ T_{s,\beta}$. Since $T_{s,\beta}$ is onto, this is exactly the
original weak equation $a_{\beta,s}(U,V)=F(V)$ for all $V\in V_y$.
The uniform bound for $T_{s,\beta}$ and the uniform coercivity give
\eqref{eq:weak-angular-inverse20}. The continuity bound of the form may
depend on $\beta$; the inverse estimate does not.
\end{proof}

\begin{corollary}[The constrained Stokes pencil]\label{cor:stokes-kernel20}
Let $K_s=\ker(d_s:V\to H_0^1(\Sigma))$. The multiplier
\[
 T_sU=W+\frac{s+2}{3-\overline s}R\omega
\]
maps $K_s$ isomorphically onto $K_t$ and satisfies
\[
 \operatorname{Re}(L_sU,L_tT_sU)
 \ge\frac{47}{192}\|\Delta_SU\|^2+\frac12y^4\|U\|^2,
 \qquad U\in K_s.
\]
Consequently the constrained weak equation on $K_s\times K_t$ has a
uniform inverse in the inherited parameter norms.
\end{corollary}

\begin{proof}
For $H=H_T+H_r\omega$, the explicit inverse is
\[
 T_s^{-1}H=H_T+\frac{3-\overline s}{s+2}H_r\omega.
\]
Since $t=1-\overline s$, direct substitution gives
\[
 d_tT_sU=d_sU,\qquad d_sT_s^{-1}H=d_tH.
\]
Hence every $H\in K_t$ equals $T_sU$ with $U=T_s^{-1}H\in K_s$,
proving surjectivity. The multiplier and its reciprocal are uniformly
bounded, so this is an isomorphism in the inherited parameter norms.
If $d_sU=0$,
the penalty term in \eqref{eq:weak-pencil20} is zero for every test.
Let $\delta\downarrow0$ in \eqref{eq:weak-angular-coercivity20}.
The transforms converge in operator norm to $T_s$, proving the estimate.
Complex Lax--Milgram on the closed space $K_s$, with the test isomorphism
$T_s$, proves the last assertion.
\end{proof}

The corollary concerns the constrained velocity equation. A full mixed
inverse additionally requires a bounded right inverse of $d_s$ and pressure
recovery. The passage from weak cone estimates to global Sobolev regularity
is carried out below using the independent straight-edge estimates
on annuli; no strong spherical-corner shift is assumed.



\section{Divergence liftings and the planar mixed input}
\label{sec:g-prelim}

All Sobolev spaces in this section are over a fixed domain. Negative-order
spaces are duals of the corresponding spaces with vanishing traces. In complex
spaces, inner products are linear in the first argument. For a bounded
connected Lipschitz domain $D$, let $\chi_D\in H^1_0(D)$ solve
$-\Delta\chi_D=1$, and set
\[
 L^2_\chi(D)=\{g\in L^2(D):(g,\chi_D)=0\},\qquad
 Q(D)=H^1_0(D)\cap L^2_0(D).
\]
The distinction between the zero-mean pressure space and the torsion-orthogonal
constraint data is necessary for the original-pressure formulation considered
below.

\begin{lemma}[A strong divergence lifting]\label{lem:g-stronglift}
Let $D$ be a bounded connected Lipschitz domain in $\mathbb R^d$, $d\ge2$. There exists
a bounded linear operator
\[
 \mathcal R_D:H^2(D)\cap H^1_0(D)\cap L^2_0(D)
 \longrightarrow H^3(D)^d\cap H^2_0(D)^d
\]
such that $\operatorname{div}\mathcal R_D h=h$.
\end{lemma}
\begin{proof}
This is \cite[Lemma~2.7]{PolygonCompanion2026}, which is stated for
bounded connected Lipschitz domains in every dimension \(d\ge2\).
Its proof constructs a bounded linear trace lifting and then applies
a support-preserving divergence inverse. In particular, the result
applies to the three-dimensional domains used here, without a
convexity or simple-connectivity assumption. For complex data we use
the complex-linear extension of this real lifting.
\end{proof}

\begin{lemma}[Mixed energy solvability]\label{lem:g-energy}
Let $D$ be bounded, connected and Lipschitz. For $F\in H^{-2}(D)^d$ and
$g\in H^{-1}(D)$ with $\langle g,\chi_D\rangle=0$, the problem
\[
 \Delta^2u+\nabla\Delta p=F,\qquad
 \Delta\operatorname{div}u=g,
 \qquad (u,p)\in H^2_0(D)^d\times Q(D)
\]
has a unique energy solution, and
\[
 \|u\|_{H^2(D)}+\|p\|_{H^1(D)}
 \le C_D\bigl(\|F\|_{H^{-2}(D)}+\|g\|_{H^{-1}(D)}\bigr).
\]
\end{lemma}
\begin{proof}
Solve $\Delta h=g$ in $H^1_0(D)$. Symmetry of the Dirichlet inverse yields
$\int_D h=-\langle g,\chi_D\rangle=0$. A support-preserving Bogovski\u\i\
operator maps $H^1_0(D)\cap L^2_0(D)$ boundedly into $H^2_0(D)^d$, so one can
lift $h$ and then solve the coercive $\Delta^2$ form on the kernel of
divergence. The residual is a bounded functional of
$\operatorname{div}v\in Q(D)$. The Riesz map for the gradient inner product
on $Q(D)$ recovers $p$, giving
\[
 (\Delta u,\Delta v)+(\nabla p,\nabla\operatorname{div}v)
 =\langle F,v\rangle\quad(v\in H^2_0(D)^d).
\]
The same lifting gives the displayed estimate. For uniqueness, the constraint
implies $\operatorname{div}u=0$ when $g=0$; testing the momentum equation by
$u$ gives $u=0$. Hence $\nabla\Delta p=0$, so $p$ is a multiple of
$\chi_D$. The mean-zero condition forces that multiple to vanish.
\end{proof}

\begin{lemma}[The full planar fourth-order Stokes shift]\label{lem:g-planarstokes}
On every bounded convex polygon $D$, the energy solution satisfies
\[
 \|u\|_{H^3(D)}+\|p\|_{H^2(D)}
 \le C_D\bigl(\|F\|_{H^{-1}(D)}+\|g\|_{L^2(D)}\bigr),
 \qquad g\in L^2_\chi(D).
\]
The scalar clamped biharmonic inverse also maps $H^{-1}(D)$ boundedly into
$H^3(D)\cap H^2_0(D)$.
\end{lemma}
\begin{proof}
The mixed estimate is \cite[Proposition~2.9]{PolygonCompanion2026}.
Its reduction from compatible nonzero constraint data to the homogeneous
constraint is proved in Proposition~2.8 there; the latter case uses
\cite[Theorem~3.6]{ChenHuangHuang2023}. The scalar clamped estimate follows from
\cite[Theorem~3.3 and Lemma~3.4]{ChenHuangHuang2023}, with \(m=2\).
\end{proof}

\begin{corollary}[Planar Lam\'e input from the companion paper]
\label{cor:planar-H320}
On a fixed bounded convex polygon \(D\), the clamped inverse of
\(P_{\lambda,\mu}\) maps \(H^{-1}(D)^2\) into
\(H^3(D)^2\cap H^2_0(D)^2\). Its norm is uniform for
\(0\le\lambda\le\Lambda<\infty\) and
\(0<\mu_0\le\mu\le\mu_1\).
\end{corollary}
\begin{proof}
This is \cite[Proposition~2.6]{PolygonCompanion2026}. The underlying
clamped planar corner pencil has no spectrum in
\(1\le\operatorname{Re}s\le2\) by Theorem~A.1 there. We use the
result for transverse edge sections; the axial-frequency and vertex
arguments remain three-dimensional and are given below.
\end{proof}

\section{The complete mixed angular inverse}
\label{sec:g-angular}

Let $\Sigma$ be the spherical section of a fixed pointed convex polyhedral
cone. It is contained, with its closure, in an open hemisphere. Put
\[
 V=H^2_0(\Sigma)^3,\quad Z=H^1_0(\Sigma),\quad
 s=a+iy,\quad \tfrac12\le a\le\tfrac32,\quad t=1-\overline s,
 \quad\kappa=(1+y^2)^{1/2}.
\]
Use Cartesian components for vector fields and write
\[
 U=U_T+U_r\omega,\quad
 d_sU=\operatorname{div}_S U_T+(s+2)U_r,
 \quad L_s=\Delta_S+s(s+1),\quad \ell_s=\Delta_S+s(s-1).
\]
The parameter norms are
\[
 \|U\|_{V_y}^2=\|\Delta_SU\|_0^2+\kappa^2\|\nabla_SU\|_0^2
                  +\kappa^4\|U\|_0^2,
 \qquad
 \|h\|_{Z_y}^2=\|\nabla_Sh\|_0^2+\kappa^2\|h\|_0^2.
\]
Their anti-duals carry the corresponding dual norms. We use the angular
coercivity established above in the following precise form: with
$N_s=\ker d_s$ and
\[
 T_sU=U_T+\frac{s+2}{3-\overline s}U_r\omega,
\]
the map $T_s:N_s\to N_t$ and its inverse are uniformly bounded in the
parameter norms, and
\[
 \operatorname{Re}(L_sU,L_tT_sU)\ge c\|U\|_{V_y}^2
 \quad(U\in N_s).
\]
No strong angular regularity on the spherical polygon is assumed.

\begin{lemma}[A parameter divergence lifting]\label{lem:g-paramlift}
For all $s$ in the closed strip, $d_s:V_y\to Z_y$ has a right inverse
$R_s$ satisfying
\[
 \|R_sh\|_{V_y}\le C_\Sigma\|h\|_{Z_y}.
\]
The analogous statement holds for $d_t$.
\end{lemma}
\begin{proof}
A gnomonic chart maps $\Sigma$ smoothly to a bounded convex polygon. At scale
$\varepsilon\asymp\kappa^{-1}$ choose a bounded-overlap cover by patches
with uniformly bounded star-shapedness constants after rescaling, and a
subordinate partition $\{\chi_j\}$ with
$|\nabla\chi_j|\le C\varepsilon^{-1}$. Such patches can be obtained by
intersecting the polygon with balls of comparable radius. At its finitely many
corners their constants depend only on the fixed corner angles. Put
$h_j=\chi_jh$ and $c_j=\int_{\Sigma_j}h_j$. The zero trace of $h$ on the
actual boundary and the support of $\chi_j$ at artificial boundaries imply
$h_j\in H^1_0(\Sigma_j)$.

Choose a smooth interior bump $\varphi_j$ with integral one. Surface dimension
two gives
\[
 \|\varphi_j\|_0\le C\varepsilon^{-1},\quad
 \|\nabla\varphi_j\|_0\le C\varepsilon^{-2},\quad
 \|D^2\varphi_j\|_0\le C\varepsilon^{-3},\quad
 |c_j|\le C\varepsilon\|h_j\|_0.
\]
Write $\Phi:D_j\to\Sigma_j$ for the inverse gnomonic chart and
$J=(\det(D\Phi^TD\Phi))^{1/2}$. The planar density
$\widehat h_j=J\,[(h_j-c_j\varphi_j)\circ\Phi]$ belongs to
$H^1_0(D_j)\cap L^2_0(D_j)$. Choose the same planar Bogovski\u\i\ operator
at the two Sobolev levels, and set
\[
 \widehat W_j=\mathcal B_j\widehat h_j,\qquad
 W_j\circ\Phi=J^{-1}D\Phi\,\widehat W_j.
\]
The surface Piola identity is
$(\operatorname{div}_S W_j)\circ\Phi=J^{-1}\operatorname{div}\widehat W_j$.
All chart coefficients and their needed derivatives are bounded on the fixed
closed polygon, with $J$ bounded away from zero. Composition and multiplication
therefore preserve $H^2_0$, as follows also by approximating
$\widehat W_j$ by compactly supported smooth fields.
Thus the tangential field $W_j\in H^2_0(\Sigma_j)^3$ obeys
\[
 \operatorname{div}_SW_j=h_j-c_j\varphi_j,
\]
\[
 \|W_j\|_0\le C\varepsilon\|h_j\|_0,\quad
 \|W_j\|_{H^1}\le C\|h_j\|_0,\quad
 \|W_j\|_{H^2}\le
 C\bigl(\|\nabla h_j\|_0+\varepsilon^{-1}\|h_j\|_0\bigr).
\]
These follow from the simultaneous $L^2_0\to H^1_0$ and
$H^1_0\cap L^2_0\to H^2_0$ bounds of the same operator in \cite{CostabelMcIntosh2010};
smooth chart coefficients add only lower-order terms of the indicated sizes.
Now set
\[
 U_j=W_j+\frac{c_j}{s+2}\varphi_j\omega.
\]
Since $|s+2|\asymp\kappa$, the radial summand has norms of orders zero, one
and two bounded respectively by
$C\varepsilon\|h_j\|_0$, $C\|h_j\|_0$ and
$C\varepsilon^{-1}\|h_j\|_0$. Moreover $d_sU_j=h_j$ exactly. Each field is
in the clamped patch space, so zero extension is valid. Bounded overlap gives
\[
 \Big\|\sum_j U_j\Big\|_{V_y}^2
 \le C\sum_j\bigl(\|\nabla h_j\|_0^2+
                         \varepsilon^{-2}\|h_j\|_0^2\bigr)
 \le C\bigl(\|\nabla h\|_0^2+\kappa^2\|h\|_0^2\bigr).
\]
Their sum is the required lifting. The reflected parameter obeys
$|t+2|\asymp\kappa$, proving the second assertion. The liftings are used for
estimates only and need not depend analytically on $s$.
\end{proof}

\begin{proposition}[Full angular isomorphism]\label{prop:g-mixedpencil}
For every $F\in V_y^*$ and $G\in Z_y^*$ there exists a unique $(U,P)\in V\times Z$
satisfying
\begin{equation}\label{eq:g-mixed-angular-system}
\begin{aligned}
 (L_sU,L_tV)-\langle\ell_sP,d_tV\rangle&=F(V)
       &&(V\in V),\\
 \ell_sd_sU&=G.
\end{aligned}
\end{equation}
Throughout the closed strip,
\begin{equation}\label{eq:g-mixed-angular-estimate}
 \|U\|_{V_y}+\|P\|_{Z_y}
 \le C_\Sigma\bigl(\|F\|_{V_y^*}+\|G\|_{Z_y^*}\bigr).
\end{equation}
On fixed unweighted angular spaces this operator and its inverse are analytic
where defined.
\end{proposition}
\begin{proof}
The Dirichlet spectral lower bound on a hemisphere and zero extension give
$\|\nabla_Sh\|_0^2\ge2\|h\|_0^2$. Since
$\operatorname{Re}[s(s-1)]\le3/4-y^2$, the operator
$\ell_s:Z_y\to Z_y^*$ is an isomorphism with uniform inverse bound.
Let $h_g=\ell_s^{-1}G$ and $U_g=R_sh_g$. On $N_s$ solve
\[
 (L_sU_0,L_tT_sW)=F(T_sW)-(L_sU_g,L_tT_sW),\qquad W\in N_s.
\]
The transformed form is bounded and coercive in $V_y$. Lax--Milgram yields
the desired bound for $U=U_0+U_g$. Its momentum residual vanishes on $N_t$.
For $k\in Z_y$ define
\[
 \langle Q,k\rangle=(L_sU,L_tR_tk)-F(R_tk).
\]
This is independent of the representative with divergence $k$, and the
parameter lifting gives $\|Q\|_{Z_y^*}\le C(\|U\|_{V_y}+\|F\|_{V_y^*})$.
Taking $P=\ell_s^{-1}Q$ proves existence and the estimate. The same
decomposition proves uniqueness.

The original pressure has degree $s-1$: if $p=r^{s-1}P$, then
$\Delta p=r^{s-3}\ell_sP$. Pairing its gradient against an angular test
gives $-\langle\ell_sP,d_tV\rangle$, which verifies the displayed pencil.
All coefficients of the resulting map
$V\times Z\to V^*\times Z^*$ are polynomial in $s$, because conjugation
in the second slot removes the conjugation in $t$. Analytic bounded-operator
inversion proves the last assertion. No angular mean-zero pressure condition
is imposed.
\end{proof}

\section{Straight-edge regularity at every axial frequency}
\label{sec:g-edges}

\begin{lemma}[Low-frequency Stokes estimate]
\label{lem:g-edge-lowfreq}
Let $W=W_\omega$, with $0<\omega<\pi$, or let $W$ be a half-plane or
$\mathbb R^2$. For $\xi\in\mathbb R$, set
\[
 L_\xi=\Delta_\perp-\xi^2,\qquad
 D_\xi=(\partial_1,\partial_2,i\xi),\qquad
 \operatorname{Div}_\xi U=\operatorname{div}_\perp U_\perp+i\xi U_3.
\]
Suppose $|\xi|\le1$ and an energy pair
$(U,p)\in H^2_0(W)^3\times H^1_0(W)$ satisfies
\[
 L_\xi^2U+D_\xi L_\xi p=F\in H^{-1}(W)^3,
 \qquad L_\xi\operatorname{Div}_\xi U=G\in L^2(W).
\]
Then $U\in H^3(W)^3$, $p\in H^2(W)$, and
\begin{equation}\label{eq:g-edge-lowfreq}
 \|U\|_{H^3(W)}+\|p\|_{H^2(W)}
 \le C_W\bigl(\|F\|_{H^{-1}(W)}+\|G\|_{L^2(W)}
                 +\|U\|_{H^2(W)}+\|p\|_{H^1(W)}\bigr),
\end{equation}
with $C_W$ independent of $\xi$. No mean condition on $p$ or $G$ is required.
\end{lemma}
\begin{proof}
Write $U=(v,w)$, $d=\operatorname{div}_\perp v$, and abbreviate
$\Delta_\perp$ and $\nabla_\perp$ by $\Delta$ and $\nabla$. In distributions,
\[
\begin{aligned}
 \Delta^2v+\nabla\Delta p
 &=F_\perp+2\xi^2\Delta v-\xi^4v+\xi^2\nabla p=:f_\xi,\\
 \Delta d
 &=G+\xi^2d-i\xi\Delta w+i\xi^3w=:g_\xi,\\
 \Delta^2w
 &=F_3+2\xi^2\Delta w-\xi^4w-i\xi\Delta p+i\xi^3p=:h_\xi.
\end{aligned}
\]
The assumed energy regularity gives
\begin{equation}\label{eq:g-edge-lowfreq-data}
 \|f_\xi\|_{H^{-1}}+\|h_\xi\|_{H^{-1}}+\|g_\xi\|_0
 \le C\bigl(\|F\|_{H^{-1}}+\|G\|_0+
                  \|U\|_{H^2}+\|p\|_{H^1}\bigr).
\end{equation}
In particular $\Delta p\in H^{-1}$ is used in its distributional sense;
no second derivative of $p$ is assumed to belong to $L^2$.

Fix a bounded-overlap cover of $W$ by bounded convex polygonal patches
$D_j$ of finitely many shapes up to rigid motions, with uniformly bounded
diameters. For the wedge, use one corner patch, translates of fixed ray
patches, and interior squares; a half-plane or the plane requires no corner
patch. Choose a subordinate partition $\{\eta_j\}$ with uniformly bounded
derivatives through order four. Each $\eta_j$ vanishes in a neighborhood of
the artificial part of $\partial D_j$, and $\sum_j\eta_j=1$ on $W$.
Choose $\varphi_j\in C_c^\infty(D_j)$ with integral one and uniformly
bounded $H^2$ norms. Set
\[
 v_j=\eta_jv,\qquad w_j=\eta_jw,\qquad
 m_j=\int_{D_j}\eta_jp,\qquad p_j=\eta_jp-m_j\varphi_j.
\]
Then $v_j\in H^2_0(D_j)^2$, $w_j\in H^2_0(D_j)$, and
$p_j\in Q(D_j)$. Their equations are
\[
\begin{aligned}
 \Delta^2v_j+\nabla\Delta p_j
 &=\eta_jf_\xi+[\Delta^2,\eta_j]v
       +[\nabla\Delta,\eta_j]p-m_j\nabla\Delta\varphi_j=:F_j,\\
 \Delta\operatorname{div}v_j
 &=\eta_jg_\xi+[\Delta\operatorname{div},\eta_j]v=:G_j,\\
 \Delta^2w_j&=\eta_jh_\xi+[\Delta^2,\eta_j]w=:H_j.
\end{aligned}
\]
The first and third right-hand sides are in $H^{-1}(D_j)$ and the second
is in $L^2(D_j)$. Indeed, the first two momentum commutators have orders
at most three in $v$ and two in $p$, respectively, and one integration
by parts gives
\[
\begin{aligned}
 \|[\Delta^2,\eta_j]v\|_{H^{-1}(D_j)}
 &\le C\|v\|_{H^2(D_j)},\\
 \|[\nabla\Delta,\eta_j]p\|_{H^{-1}(D_j)}
 &\le C\|p\|_{H^1(D_j)},\\
 \|[\Delta\operatorname{div},\eta_j]v\|_{L^2(D_j)}
 &\le C\|v\|_{H^2(D_j)}.
\end{aligned}
\]
The scalar commutator has the same bound as the first one. Furthermore,
$|m_j|\le C\|p\|_{L^2(D_j)}$ and
$\|m_j\nabla\Delta\varphi_j\|_{H^{-1}(D_j)}\le C|m_j|$.
Since $\operatorname{div}v_j\in H^1_0(D_j)$, the constraint is automatically
compatible:
\[
 (G_j,\chi_{D_j})=-\int_{D_j}\operatorname{div}v_j=0.
\]

For any $f\in H^{-1}(W)$,
\begin{equation}\label{eq:g-edge-patch-duality}
 \sum_j\|\eta_j f\|_{H^{-1}(D_j)}^2
 \le C\|f\|_{H^{-1}(W)}^2.
\end{equation}
To see this, extend tests by zero and use the bounded map
$(\phi_j)_j\mapsto\sum_j\eta_j\phi_j$ from the Hilbert sum of
$H^1_0(D_j)$ into $H^1_0(W)$, then take its adjoint. The remaining local
terms square-sum by bounded overlap and the uniform cutoff and bump bounds.
Thus~\eqref{eq:g-edge-lowfreq-data} controls
\[
 \sum_j\bigl(\|F_j\|_{H^{-1}}^2+\|G_j\|_0^2+
                 \|H_j\|_{H^{-1}}^2+|m_j|^2\bigr).
\]
Apply Lemma~\ref{lem:g-planarstokes} to the already given energy pair
$(v_j,p_j)$ and its scalar assertion to $w_j$. Its constants are uniform
in $j$ because the patches have finitely many shapes. Reconstruct
$v=\sum_jv_j$, $w=\sum_jw_j$, and
$p=\sum_j(p_j+m_j\varphi_j)$. After obtaining the local regularity, these
fields may be extended across artificial patch boundaries: each vanishes
in a neighborhood of those boundaries. No extension across the physical
boundary of $W$ is used. Bounded overlap now proves
\eqref{eq:g-edge-lowfreq}. Every right-hand side was estimated using only
the original energy norms, so no regularity bootstrap was assumed.
\end{proof}

\begin{proposition}[The fourth-order Stokes edge estimate]\label{prop:g-edge}
Let $W_\omega=\{r>0,\ 0<\theta<\omega\}$, $0<\omega<\pi$, and
$\mathcal W=W_\omega\times\mathbb R$. Suppose an energy pair
$(u,p)\in H^2_0(\mathcal W)^3\times H^1_0(\mathcal W)$ satisfies
\[
 \Delta^2u+\nabla\Delta p=F\in H^{-1}(\mathcal W)^3,
 \qquad \Delta\operatorname{div}u=G\in L^2(\mathcal W).
\]
Then $u\in H^3$, $p\in H^2$, and
\begin{equation}\label{eq:g-edge-estimate}
 \|u\|_{H^3}+\|p\|_{H^2}
 \le C_\omega\bigl(\|F\|_{H^{-1}}+\|G\|_0+
                              \|u\|_{H^2}+\|p\|_{H^1}\bigr).
\end{equation}
The same assertion holds for a half-plane or $\mathbb R^2$ as transverse
section.
\end{proposition}
\begin{proof}
After Fourier transformation along the edge, write
\[
 L_\xi=\Delta_\perp-\xi^2,\qquad
 D_\xi=(\partial_1,\partial_2,i\xi),\qquad
 \operatorname{Div}_\xi U=\operatorname{div}_\perp U_\perp+i\xi U_3.
\]
The equations become
\begin{equation}\label{eq:g-edge-fourier-reduction}
 L_\xi^2U+D_\xi L_\xi p=F,\qquad
 L_\xi\operatorname{Div}_\xi U=G.
\end{equation}

\emph{Unit frequency.}
For $\xi=\sigma\in\{-1,1\}$, the map
$\operatorname{Div}_\sigma:H^2_0(W_\omega)^3\to H^1_0(W_\omega)$ has a
bounded right inverse. Indeed cover the wedge by uniformly bounded convex
polygonal patches of finitely many shapes: one corner patch, ray patches,
interior squares and finitely many transition patches. Choose a subordinate
smooth partition with uniformly bounded derivatives. If $h_j=\chi_jh$ and
$c_j=\int h_j$, lift $h_j-c_j\varphi_j$ by planar Bogovski\u\i\ and put
$c_j\varphi_j/(i\sigma)$ in the third component. Zero extension and
bounded-overlap summation give the right inverse. Here each bump has integral
one and uniformly bounded Sobolev norms. All local means are lifted; no total
mean is required on the infinite wedge.

The form $\|(\Delta_\perp-1)U\|_0^2$ is coercive in $H^2_0$, while
$-(\Delta_\perp-1)$ is coercive in $H^1_0$. The preceding right inverse
therefore gives unit-frequency mixed energy solvability and
\[
 \|U\|_{H^2}+\|p\|_{H^1}
 \le C_\omega\bigl(\|F\|_{H^{-2}}+\|G\|_{H^{-1}}\bigr).
\]
Applying Lemma~\ref{lem:g-edge-lowfreq} at $\xi=\sigma\in\{-1,1\}$ and
using the preceding mixed energy estimate, with
$\|F\|_{H^{-2}}\le C\|F\|_{H^{-1}}$ and
$\|G\|_{H^{-1}}\le\|G\|_0$, yields
\begin{equation}\label{eq:g-edge-unit-estimate}
 \|U\|_{H^3}+\|p\|_{H^2}
 \le C_\omega\bigl(\|F\|_{H^{-1}}+\|G\|_0\bigr)
 \qquad(|\xi|=1).
\end{equation}

\emph{Scaling and low frequencies.}
For $\tau=|\xi|>0$ set
\begin{equation}\label{eq:g-edge-fourier-rescaling}
 \widetilde U(X)=U(X/\tau),\qquad
 \widetilde p(X)=\tau^{-1}p(X/\tau),\qquad
 \widetilde F=\tau^{-4}F(X/\tau),\qquad
 \widetilde G=\tau^{-3}G(X/\tau).
\end{equation}
These satisfy the unit-frequency system. Write
$\|v\|_{k,\tau}^2=\sum_{j=0}^k\tau^{2(k-j)}|v|_j^2$, with
$|v|_0=\|v\|_0$. The norm $\|\cdot\|_{-1,\tau}$ is dual to
$(\|\nabla\phi\|_0^2+\tau^2\|\phi\|_0^2)^{1/2}$ on
$H^1_0(W_\omega)$. All four norms in
\eqref{eq:g-edge-unit-estimate} scale by $\tau^{-2}$, because the
transverse dimension is two. Consequently,
\begin{equation}\label{eq:g-edge-scaled-estimate}
 \|U\|_{3,\tau}+\|p\|_{2,\tau}
 \le C_\omega\bigl(\|F\|_{-1,\tau}+\|G\|_0\bigr),\qquad\tau>0.
\end{equation}
For $|\xi|\ge1$ these norms are equivalent to those with parameter
$\kappa=(1+\xi^2)^{1/2}$. For $|\xi|\le1$, Lemma~\ref{lem:g-edge-lowfreq} applies to the given energy
pair. On this interval the parameter norms and the ordinary Sobolev norms,
including their dual norms, are uniformly equivalent. Hence
\[
 \|U\|_{3,\kappa}+\|p\|_{2,\kappa}
 \le C_\omega\bigl(\|F\|_{-1,\kappa}+\|G\|_0+
                 \|U\|_{2,\kappa}+\|p\|_{1,\kappa}\bigr).
\]
Together with the high-frequency estimate this is uniform for all
$\xi\in\mathbb R$. Fourier--Plancherel and the dual direct-integral identity
then give~\eqref{eq:g-edge-estimate}. The negative norm in the Fourier
integral is always dual to the inhomogeneous parameter norm with
$\kappa=(1+\xi^2)^{1/2}$. Half-plane and interior patches are included in
Lemma~\ref{lem:g-edge-lowfreq} and the same unit-frequency argument.
\end{proof}

\section{A joint Mellin shift and the global fourth-order Stokes theorem}
\label{sec:g-global}

\begin{lemma}[Cylinder inversion and a pole-free Mellin shift]
\label{lem:g-mellin-inversion}
Let $\Sigma,V,Z$ and the mixed angular pencil be as in
Proposition~\ref{prop:g-mixedpencil}. Put
\[
 a_0=\tfrac12,\qquad a_1=\tfrac32,\qquad
 \mathcal C=\mathbb R\times\Sigma,\qquad E=V\times Z,
\]
and define the cylindrical spaces
\[
 \mathcal X=H^2_0(\mathcal C)^3\times H^1_0(\mathcal C),
 \qquad
 \mathcal X^*=(H^{-2}_{\rm cyl})^3\times H^{-1}_{\rm cyl}.
\]
Here $H^m_0(\mathcal C)$ is the closure of $C_c^\infty(\mathcal C)$
in the full cylindrical Sobolev norm, and $H^{-m}_{\rm cyl}$ is its
antidual. Thus the velocity is clamped and the pressure has zero
Dirichlet trace on the lateral boundary; no angular mean constraint is
imposed.

Denote the operator in \eqref{eq:g-mixed-angular-system} by
$\mathcal A(s):E\to E^*$. For every $a\in[a_0,a_1]$, the Fourier
multiplier
\[
 \widehat{\mathcal B_a Y}(y)=
 \mathcal A(a+iy)\widehat Y(y)
\]
defines an isomorphism $\mathcal B_a:\mathcal X\to\mathcal X^*$, with
\begin{equation}\label{eq:g-cylinder-line-inverse}
 \|\mathcal B_a^{-1}\mathcal D\|_{\mathcal X}
 \le C_\Sigma\|\mathcal D\|_{\mathcal X^*}
 \qquad(a_0\le a\le a_1).
\end{equation}
The angular inverse is holomorphic on an open set containing the closed
strip $a_0\le\operatorname{Re}s\le a_1$, and its fixed-space norm
$E^*\to E$ is uniformly bounded on that strip.

Suppose that $\mathcal D_1\in\mathcal X^*$ has support in
$\{\tau\le T\}$. Set
\[
 \mathcal D_a=e^{(a_1-a)\tau}\mathcal D_1,\qquad
 Y_a=\mathcal B_a^{-1}\mathcal D_a.
\]
Then
\begin{equation}\label{eq:g-mellin-line-identification}
 e^{a\tau}Y_a=e^{b\tau}Y_b
 \quad\hbox{in }L^2_{\rm loc}(\mathbb R;E),
 \qquad a,b\in[a_0,a_1].
\end{equation}
In particular, if an already given energy-line pair $Y_0\in\mathcal X$
solves
$\mathcal B_{a_0}Y_0=e^{(a_1-a_0)\tau}\mathcal D_1$, then
\[
 e^{-(a_1-a_0)\tau}Y_0\in\mathcal X,\qquad
 \|e^{-(a_1-a_0)\tau}Y_0\|_{\mathcal X}
 \le C_\Sigma\|\mathcal D_1\|_{\mathcal X^*}.
\]
\end{lemma}

\begin{proof}
Use $\widehat f(y)=\int_{\mathbb R}e^{-iy\tau}f(\tau)\,d\tau$.
Plancherel's theorem and the definitions of $V_y,Z_y$ give
\begin{align}
 \|(U,P)\|_{\mathcal X}^2
 &\asymp\int_{\mathbb R}
 \bigl(\|\widehat U(y)\|_{V_y}^2+
       \|\widehat P(y)\|_{Z_y}^2\bigr)\,dy,
 \label{eq:g-cylindrical-plancherel}\\
 \|(\mathcal F,\mathcal G)\|_{\mathcal X^*}^2
 &\asymp\int_{\mathbb R}
 \bigl(\|\widehat{\mathcal F}(y)\|_{V_y^*}^2+
       \|\widehat{\mathcal G}(y)\|_{Z_y^*}^2\bigr)\,dy.
 \label{eq:g-cylindrical-duality}
\end{align}
For the first identity, the velocity terms are the squared angular
second derivatives, the mixed radial--angular derivatives, and the
radial derivatives of orders at most two. The angular $\Delta_S$ norm
is equivalent to the $H^2_0(\Sigma)$ norm by zero extension to the
hemisphere and the sphere's integration-by-parts identities. The
pressure identity is the corresponding first-order formula.
Tensor products of compactly supported smooth functions are dense in
the cylindrical spaces. Hilbert-space duality therefore gives the
second identity, including the negative-order data spaces.

The product rule, $|s|\le C(1+|y|)$, and the definitions of the parameter
norms imply, uniformly on the closed strip,
\[
 \|L_sU\|_0+\|d_sU\|_{Z_y}\le C_\Sigma\|U\|_{V_y},
 \qquad
 \|\ell_sh\|_{Z_y^*}\le C_\Sigma\|h\|_{Z_y},
\]
and the analogous bounds hold for $t=1-\overline s$.
Thus $\mathcal A(s)$ is uniformly bounded from
$V_y\times Z_y$ to $V_y^*\times Z_y^*$.
Proposition~\ref{prop:g-mixedpencil} supplies the inverse bound.
The two direct-integral identities now prove the asserted line
isomorphism and \eqref{eq:g-cylinder-line-inverse}.

On the fixed spaces $E,E^*$, all coefficients of $\mathcal A(s)$ are
polynomial in $s$: conjugation in the second slot removes the
conjugation in $t$. Its invertibility set is open by the bounded
inverse theorem and contains every point of the closed strip.
Analytic operator inversion gives a holomorphic inverse on that open
set. Moreover, the parameter solution norms dominate the fixed
$E$ norm, while the parameter dual norms are bounded by the fixed
$E^*$ norm. Hence
\[
 \sup_{a_0\le\operatorname{Re}s\le a_1}
 \|\mathcal A(s)^{-1}\|_{\mathcal L(E^*,E)}<\infty.
\]
No bound outside the closed strip is needed below: every finite
closed rectangle in it has an open neighborhood on which the inverse
is holomorphic.

It remains to identify the solutions on different lines. All
multiplications by $e^{(a_1-a)\tau}$ below are understood on data
supported in $\{\tau\le T+1\}$, using a fixed smooth upper cutoff.
Their derivatives through order two are then uniformly bounded for
$a\in[a_0,a_1]$, so they are bounded multipliers on $\mathcal X^*$.
Choose smooth lower cutoffs converging to one, with uniformly bounded
derivatives, and then convolve in $\tau$ with a compactly supported
approximate identity. These operations give
\[
 \mathcal D_{1,n}\in C_c^\infty(\mathbb R;E^*),\qquad
 \operatorname{supp}\mathcal D_{1,n}\subset\{\tau\le T+1\},
 \qquad
 \mathcal D_{1,n}\longrightarrow\mathcal D_1
       \quad\hbox{in }\mathcal X^*.
\]
Indeed
\[
 \mathcal X^*\hookrightarrow
 H^{-2}(\mathbb R;V^*)\times H^{-1}(\mathbb R;Z^*),
\]
by duality from
$H^2(\mathbb R;V)\hookrightarrow H^2_0(\mathcal C)^3$
and
$H^1(\mathbb R;Z)\hookrightarrow H^1_0(\mathcal C)$.
Radial convolution therefore produces smooth $E^*$-valued functions;
angular smoothing is unnecessary. Strong convergence of the lower
cutoffs and convolution in $\mathcal X^*$ follows first for compactly
supported smooth data and then by density and their uniform multiplier
bounds.

Put
\[
 \mathcal D_{a,n}=e^{(a_1-a)\tau}\mathcal D_{1,n},\qquad
 Y_{a,n}=\mathcal B_a^{-1}\mathcal D_{a,n},\qquad
 H_n=e^{a_1\tau}\mathcal D_{1,n}.
\]
For every fixed $a$, both $\mathcal D_{a,n}\to\mathcal D_a$ in
$\mathcal X^*$ and $Y_{a,n}\to Y_a$ in $\mathcal X$.
The $E^*$-valued bilateral Laplace transform
\[
 \mathcal L H_n(s)=\int_{\mathbb R}e^{-s\tau}H_n(\tau)\,d\tau
\]
is entire. Repeated integration by parts gives, for every integer
$N\ge0$,
\[
 \|\mathcal L H_n(a+iy)\|_{E^*}
 \le C_{n,N}(1+|y|)^{-N},
 \qquad a_0\le a\le a_1.
\]
The same decay holds for
$\mathcal A(a+iy)^{-1}\mathcal L H_n(a+iy)$ in $E$.
Fourier inversion thus gives the absolutely convergent $E$-valued
formula
\[
 e^{a\tau}Y_{a,n}(\tau)=
 \frac1{2\pi i}\int_{\operatorname{Re}s=a}
 e^{s\tau}\mathcal A(s)^{-1}\mathcal L H_n(s)\,ds .
\]
Cauchy's theorem on a rectangle between any two lines
$a,b\in[a_0,a_1]$ identifies these integrals. The horizontal integrals
tend to zero by the rapid decay and the uniform fixed-space inverse
bound. Consequently
$e^{a\tau}Y_{a,n}=e^{b\tau}Y_{b,n}$.
On every bounded $\tau$ interval the exponential multipliers are
bounded, and $\mathcal X\hookrightarrow L^2(\mathbb R;E)$.
Passing to the limit proves \eqref{eq:g-mellin-line-identification}.
Uniqueness of the line inverse identifies any given energy-line
solution with $Y_{a_0}$. The last estimate is then
\eqref{eq:g-cylinder-line-inverse} at $a=a_1$.
\end{proof}

\begin{lemma}[Weighted vertex estimate]\label{lem:g-vertexweights}
Let $K$ be a pointed convex polyhedral cone, and let an energy pair
$(v,q)\in H^2_0(K)^3\times H^1_0(K)$ have bounded outer radial support and
satisfy
\[
 \Delta^2v+\nabla\Delta q=F\in H^{-1}(K)^3,\qquad
 \Delta\operatorname{div}v=G\in L^2(K).
\]
Then, with $r=|x|$,
\[
 \|r^{-1}D^2v\|_0+\|r^{-2}Dv\|_0+\|r^{-3}v\|_0
 +\|r^{-1}Dq\|_0+\|r^{-2}q\|_0
 \le C\bigl(\|F\|_{H^{-1}(K)}+\|G\|_0\bigr).
\]
The constant may depend on the fixed cone and the outer support bound.
\end{lemma}
\begin{proof}
Set $\tau=\log r$ and $\mathcal C=\mathbb R\times\Sigma$.
For smooth compactly supported fields, $dx=r^3\,d\tau\,d\omega$ and
$\partial_\tau=r\partial_r$ give
\begin{align}
 \|r^{-a}v\|_{H^2(\mathcal C)}^2
 &\asymp\sum_{j=0}^{2}\int_K
       r^{2j-2a-3}|D^jv|^2\,dx,
 \label{eq:g-cylindrical-velocity}\\
 \|r^{1-a}q\|_{H^1(\mathcal C)}^2
 &\asymp\sum_{j=0}^{1}\int_K
       r^{2j-2a-1}|D^jq|^2\,dx .
 \label{eq:g-cylindrical-pressure}
\end{align}
Here the expressions on the left denote cylindrical pullbacks.
The constants are uniform for $a\in[1/2,3/2]$, and the identities extend
to the associated weighted spaces by density.

Let $a_0=1/2$. Hardy's inequality applied to $Dv$ and $q$, after
zero extension from the cone, gives $Dv/r,q/r\in L^2(K)$.
Angular Poincar\'e applied to $v$ then gives $v/r^2\in L^2(K)$.
It follows that the already given energy pair determines
\[
 Y_0=(r^{-1/2}v,r^{1/2}q)\in
 H^2_0(\mathcal C)^3\times H^1_0(\mathcal C)=\mathcal X.
\]
Its lateral boundary conditions follow by transforming the compactly
supported smooth approximations in the physical energy spaces; no
additional trace condition is imposed at the vertex.

Put $a_1=3/2$. The cylindrical data on a general line are
\begin{equation}\label{eq:g-mellin-data-conversion}
 F_a=r^{4-a}F,\qquad G_a=r^{3-a}G.
\end{equation}
For the distributional force at the target line, this means
\[
 \langle F_1,V\rangle
 =\langle F,\eta(r)r^{-1/2}V(\log r,\omega)\rangle ,
\]
where $\eta$ is a fixed smooth outer radial cutoff equal to one near
the support of $F$. If the common outer support radius is $R_0$, then
\[
 \|\eta r^{-1/2}V\|_{H^1(K)}
 \le C_{R_0}\|V\|_{H^1(\mathcal C)}.
\]
Indeed the differentiated weight is $r^{-3/2}$, whose square cancels
the cylindrical Jacobian $r^3$; the undifferentiated and cutoff terms
are bounded using the outer support. Density identifies the physical
test as an element of $H^1_0(K)$. Hence
\[
 \|F_1\|_{H^{-1}_{\rm cyl}}\le C_{R_0}\|F\|_{H^{-1}(K)},
 \qquad
 \|G_1\|_{L^2(\mathcal C)}=\|G\|_{L^2(K)}.
\]
In particular $\mathcal D_1=(F_1,G_1)\in\mathcal X^*$, with
\[
 \|\mathcal D_1\|_{\mathcal X^*}
 \le C_{R_0}\bigl(\|F\|_{H^{-1}(K)}+\|G\|_0\bigr).
\]
Both components have a common upper support bound in $\tau$, and
$\mathcal D_0=e^{(a_1-a_0)\tau}\mathcal D_1$.

The homogeneous identities defining the angular pencil, with $s$
replaced by $a+\partial_\tau$, give the conjugated equation
$\mathcal B_{a_0}Y_0=\mathcal D_0$. This identity holds first for
smooth fields and then distributionally in the energy spaces.
Lemma~\ref{lem:g-mellin-inversion} therefore identifies the target-line
solution as
\[
 Y_1=e^{-\tau}Y_0=(r^{-3/2}v,r^{-1/2}q)\in\mathcal X,
 \qquad
 \|Y_1\|_{\mathcal X}\le C_\Sigma\|\mathcal D_1\|_{\mathcal X^*}.
\]
Finally, \eqref{eq:g-cylindrical-velocity}--%
\eqref{eq:g-cylindrical-pressure} at $a=3/2$ identify this norm with
the five weighted norms in the statement. This proves the estimate
without assuming target regularity of the energy pair.
\end{proof}

\begin{theorem}[Global fourth-order Stokes regularity]\label{thm:g-stokes}
Let $\Omega\subset\mathbb R^3$ be a fixed bounded convex polyhedron. For
$F\in H^{-1}(\Omega)^3$ and $g\in L^2_\chi(\Omega)$, its energy solution
satisfies
\[
 \|u\|_{H^3(\Omega)}+\|p\|_{H^2(\Omega)}
 \le C_\Omega\bigl(\|F\|_{H^{-1}(\Omega)}+\|g\|_{L^2(\Omega)}\bigr).
\]
Here $u\in H^2_0(\Omega)^3$ and $p\in H^1_0(\Omega)\cap L^2_0(\Omega)$.
\end{theorem}
\begin{proof}
Near a vertex choose a common smooth radial cutoff and write
$v=\chi u$, $q=\chi p$. They extend to the infinite tangent cone in their
energy spaces. The localized momentum force belongs to $H^{-1}$ because its
commutators have at most three derivatives of $u$ and two of $p$. The localized
constraint belongs to $L^2$, since
\[
 \begin{split}
 \Delta\operatorname{div}(\chi u)-\chi\Delta\operatorname{div}u
 ={}&2\nabla\chi\cdot\nabla\operatorname{div}u
 +(\Delta\chi)\operatorname{div}u+\nabla\chi\cdot\Delta u\\
 &+2\sum_{i,j}\partial_{ij}\chi\,\partial_i u_j
 +\nabla\Delta\chi\cdot u.
 \end{split}
\]
Only second derivatives of the energy velocity occur here. The localized
pressure need not have zero mean. Lemma~\ref{lem:g-vertexweights} applies to
this pair with bounds by the original data and energy norms.

On fixed nested cone annuli $A\Subset_r A^*$, which are separated from the
vertex, Proposition~\ref{prop:g-edge}, its flat-face and interior versions,
and a finite partition give
\[
 \|v\|_{H^3(A)}+\|q\|_{H^2(A)}
 \le C\bigl(\|F\|_{H^{-1}(A^*)}+\|G\|_{L^2(A^*)}
                  +\|v\|_{H^2(A^*)}+\|q\|_{H^1(A^*)}\bigr).
\]
All cutoffs vanish near artificial radial boundaries; no boundary condition
is imposed there. For $A_h=hA$, set
\[
 U(X)=v(hX),\qquad P(X)=h q(hX),\qquad
 F_h=h^4F(hX),\qquad G_h=h^3G(hX).
\]
In dimension three, the $D^3U$, $D^2P$, $F_h$ in $H^{-1}$, and $G_h$ in
$L^2$ norms each carry the factor $h^{3/2}$, provided the physical force uses
the dual norm of
$(\|\nabla\phi\|_0^2+h^{-2}\|\phi\|_0^2)^{1/2}$. Dividing by this factor
shows that the energy remainders are exactly
\[
 h^{-1}\|D^2v\|_0+h^{-2}\|Dv\|_0+h^{-3}\|v\|_0
 +h^{-1}\|Dq\|_0+h^{-2}\|q\|_0.
\]
Their squared sums over bounded-overlap dyadic annuli are controlled by
Lemma~\ref{lem:g-vertexweights}.

For completeness, the force square sum is bounded by
\begin{equation}\label{eq:g-force-dyadic-dual}
 \sum_j\|F\|_{H^{-1}_{1/h_j}(A_{h_j}^*)}^2
 \le C\|F\|_{H^{-1}(K)}^2.
\end{equation}
Indeed the sum of zero extensions maps the Hilbert sum of the weighted
$H^1_0(A_{h_j}^*)$ test spaces boundedly into $H^1_0(K)$, by finite overlap
and the bounded outer scale. Take its Hilbert-space adjoint. This treats
distributional forces and does not sum infinitely many copies of one global
bound. The $G$ terms sum directly in $L^2$.

We obtain $H^3\times H^2$ near every vertex. A finite cover treats the
remaining straight edges, faces and interior. This yields the stated global
estimate with an additional $H^2\times H^1$ remainder, which is removed by
Lemma~\ref{lem:g-energy}. All third velocity derivatives and second pressure
derivatives belong to $L^2$ on the annular exhaustion; the omitted vertex is
a boundary point, so this also gives the required distributional Sobolev
membership on the domain.
\end{proof}

\section{Uniform fourth-order Lam\'e regularity}
\label{sec:g-lameuniform}

\begin{lemma}[Lam\'e edge estimate on bounded material intervals]
\label{lem:g-lame-edge-bounded}
Let $W$ be a convex wedge, a half-plane, or $\mathbb R^2$, and put
$\mathcal W=W\times\mathbb R$. For each $B\ge1$ there is a constant
$C_{W,B}$ such that, whenever $1\le\beta\le B$ and
$u\in H^2_0(\mathcal W)^3$ satisfies
\[
 (\Delta^2+\beta\nabla\Delta\operatorname{div})u
       =F\in H^{-1}(\mathcal W)^3,
\]
we have $u\in H^3(\mathcal W)^3$ and
\begin{equation}\label{eq:g-lame-edge-bounded}
 \|u\|_{H^3(\mathcal W)}
 \le C_{W,B}\bigl(\|F\|_{H^{-1}(\mathcal W)}
                         +\|u\|_{H^2(\mathcal W)}\bigr).
\end{equation}
\end{lemma}
\begin{proof}
Fourier transformation in the axial variable gives
\[
 \mathcal B_{\beta,\xi}U
 :=L_\xi^2U+\beta D_\xi L_\xi\operatorname{Div}_\xi U=F.
\]
First let $|\xi|\le1$, write $U=(v,w)$ and
$d=\operatorname{div}_\perp v$, and again abbreviate transverse derivatives.
The equations are exactly
\[
\begin{aligned}
 (\Delta^2+\beta\nabla\Delta\operatorname{div})v
 ={}&F_\perp+2\xi^2\Delta v-\xi^4v
       +\beta\xi^2\nabla d
       -i\beta\xi\nabla\Delta w+i\beta\xi^3\nabla w=:f_{\beta,\xi},\\
 \Delta^2w
 ={}&F_3+(2+\beta)\xi^2\Delta w-(1+\beta)\xi^4w
       -i\beta\xi\Delta d+i\beta\xi^3d=:h_{\beta,\xi}.
\end{aligned}
\]
In particular,
\[
 \|f_{\beta,\xi}\|_{H^{-1}(W)}
       +\|h_{\beta,\xi}\|_{H^{-1}(W)}
 \le C_B\bigl(\|F\|_{H^{-1}(W)}+\|U\|_{H^2(W)}\bigr).
\]
The third-order terms $\nabla\Delta w$ and $\Delta d$ are distributions
in $H^{-1}$ controlled by $\|U\|_{H^2}$; neither component is first
assumed to have $H^3$ regularity.

Use the patches and cutoffs of Lemma~\ref{lem:g-edge-lowfreq}. For
$\mathcal P^\perp_\beta=\Delta^2+\beta\nabla\Delta\operatorname{div}$,
\[
\begin{aligned}
 \mathcal P^\perp_\beta(\eta_jv)
 &=\eta_jf_{\beta,\xi}+[\mathcal P^\perp_\beta,\eta_j]v,\\
 \Delta^2(\eta_jw)
 &=\eta_jh_{\beta,\xi}+[\Delta^2,\eta_j]w.
\end{aligned}
\]
Their commutators have order at most three and satisfy
\[
 \|[\mathcal P^\perp_\beta,\eta_j]v\|_{H^{-1}(D_j)}
       +\|[\Delta^2,\eta_j]w\|_{H^{-1}(D_j)}
 \le C_B\|U\|_{H^2(D_j)}.
\]
Corollary~\ref{cor:planar-H320} applies with planar Lam\'e parameters
$\mu=1$ and $\lambda=\beta-1\in[0,B-1]$; its constants are uniform for
these parameters and the finitely many patch shapes. Apply it to
$\eta_jv$, and the scalar part of Lemma~\ref{lem:g-planarstokes} to
$\eta_jw$. The dual localization estimate
\eqref{eq:g-edge-patch-duality} and bounded overlap give
\begin{equation}\label{eq:g-lame-lowfreq}
 \|U\|_{H^3(W)}
 \le C_{W,B}\bigl(\|F\|_{H^{-1}(W)}+\|U\|_{H^2(W)}\bigr),
 \qquad |\xi|\le1.
\end{equation}
As in Lemma~\ref{lem:g-edge-lowfreq}, only artificial boundaries are crossed
in reconstructing the local fields.

For $\xi=\sigma\in\{-1,1\}$, testing the energy equation with $U$ gives
\[
 \|L_\sigma U\|_0^2
  +\beta\|D_\sigma\operatorname{Div}_\sigma U\|_0^2
  =\operatorname{Re}\langle F,U\rangle.
\]
The first term equals
$\|\Delta_\perp U\|_0^2+2\|\nabla_\perp U\|_0^2+\|U\|_0^2$
and controls $\|U\|_{H^2(W)}^2$. Hence
$\|U\|_{H^2}\le C_W\|F\|_{H^{-2}}$, uniformly in $\beta\ge1$.
Together with~\eqref{eq:g-lame-lowfreq}, this proves
\[
 \|U\|_{H^3(W)}\le C_{W,B}\|F\|_{H^{-1}(W)},
 \qquad |\xi|=1.
\]
For $\tau=|\xi|>0$, the rescaling
\[
 \widetilde U(X)=U(X/\tau),\qquad
 \widetilde F(X)=\tau^{-4}F(X/\tau)
\]
reduces the equation to frequency $\xi/\tau$ with the same $\beta$.
The two norms scale by the common factor $\tau^{-2}$, giving
\[
 \|U\|_{3,\tau}\le C_{W,B}\|F\|_{-1,\tau}.
\]
For $|\xi|\ge1$ these norms are uniformly equivalent to the norms with
$\kappa=(1+\xi^2)^{1/2}$. For $|\xi|\le1$, use
\eqref{eq:g-lame-lowfreq} and the corresponding equivalence with ordinary
Sobolev norms. Thus, for every real $\xi$,
\[
 \|U\|_{3,\kappa}
 \le C_{W,B}\bigl(\|F\|_{-1,\kappa}+\|U\|_{2,\kappa}\bigr).
\]
The direct-integral identities for the full Sobolev and dual norms now
prove~\eqref{eq:g-lame-edge-bounded}. Uniformity is asserted only on
$1\le\beta\le B$; no localization estimate uniform as $B\to\infty$ is used.
\end{proof}

\begin{theorem}\label{thm:g-lameuniform}
On a fixed bounded convex polyhedron $\Omega$, let $u\in H^2_0(\Omega)^3$
solve
\[
 \mu\Delta^2u+(\lambda+\mu)\nabla\Delta\operatorname{div}u=F,
 \qquad F\in H^{-1}(\Omega)^3.
\]
For all $\lambda\ge0$ and $\mu>0$,
\[
 \mu\|u\|_{H^3(\Omega)}
 +(\lambda+\mu)\|\operatorname{div}u\|_{H^2(\Omega)}
 \le C_\Omega\|F\|_{H^{-1}(\Omega)}.
\]
\end{theorem}
\begin{proof}
First fix a finite $\beta=(\lambda+\mu)/\mu\ge1$. The angular
coercivity of Theorem~\ref{thm:weak-angular20} gives its weak inverse on the entire
closed strip $1/2\le\operatorname{Re}s\le3/2$. The proofs of
Lemmas~\ref{lem:g-mellin-inversion} and~\ref{lem:g-vertexweights}, with the
pressure and constraint components omitted and this vector pencil used in the
contour integral, therefore give
$D^2v/r$, $Dv/r^2$ and $v/r^3$ in $L^2$ for a localized energy solution.

At straight edges, Lemma~\ref{lem:g-lame-edge-bounded} gives the local
$H^3$ estimate for each fixed finite $\beta$ and, more generally, uniformly
on every fixed interval $1\le\beta\le B$. Its half-plane and interior
versions treat faces and interior patches. Combining these estimates with
the vertex weights in the preceding paragraph and the dyadic annular
argument of Theorem~\ref{thm:g-stokes} proves the
$H^{-1}\to H^3\cap H^2_0$ shift for each finite $\beta$.

To obtain a constant uniform as $\beta\to\infty$, set $\gamma=\beta^{-1}$
and use the fixed strong spaces
\[
 \mathcal X=(H^3\cap H^2_0)^3\times(H^2\cap H^1_0\cap L^2_0),\qquad
 \mathcal Z=H^{-1}(\Omega)^3\times L^2_\chi(\Omega).
\]
By Theorem~\ref{thm:g-stokes},
\[
 S_0(u,p)=(\Delta^2u+\nabla\Delta p,\ \Delta\operatorname{div}u)
 :\mathcal X\longrightarrow\mathcal Z
\]
is a bounded isomorphism. The map $B(u,p)=(0,\Delta p)$ is bounded into
$\mathcal Z$, because
$(\Delta p,\chi_\Omega)=-(p,1)=0$. Hence
\begin{equation}\label{eq:g-strong-space-perturbation}
 S_\gamma=S_0-\gamma B
\end{equation}
has a uniformly bounded inverse for sufficiently small $\gamma\ge0$. Applied to data $(F/\mu,0)$, its second equation and
zero traces imply
\[
 \Delta(\operatorname{div}u-\gamma p)=0,
 \qquad \operatorname{div}u-\gamma p\in H^1_0(\Omega),
\]
so $p=\beta\operatorname{div}u$. Energy uniqueness identifies this strong
solution with the given Lam\'e solution. It gives
$\|u\|_{H^3}+\|\beta\operatorname{div}u\|_{H^2}
\le C\mu^{-1}\|F\|_{H^{-1}}$ on the large-$\beta$ range.

The remaining $\beta$ values form a compact interval. Each finite-$\beta$
operator is an isomorphism by the first part of the proof and energy
solvability. These operators depend continuously on $\beta$ in
$\mathcal L(H^3\cap H^2_0,H^{-1})$; bounded-operator inversion and compactness
give a uniform inverse bound on that interval. There $\beta$ itself is bounded,
so the divergence term is controlled by the $H^3$ norm. Multiplying the two
range estimates by $\mu$ proves the result. Thus no uniform bound for
unbounded-$\beta$ localization commutators is assumed.
\end{proof}



\section{Uniform classical elasticity on convex polyhedra}
\label{sec:est-classical}

All Sobolev spaces in this section are complexified when Fourier or Mellin
transforms are used. Inner products are linear in their first argument.
Constants may depend on the fixed domain, but not on the material parameters.

\begin{theorem}\label{thm:est-classical}
Let $\Omega\subset\mathbb R^3$ be a fixed bounded convex polyhedron.
For $\lambda\ge0$, $\mu>0$ and $f\in L^2(\Omega)^3$, let
$u_0\in H_0^1(\Omega)^3$ solve
\[
 -\mu\Delta u_0-(\lambda+\mu)\nabla\operatorname{div}u_0=f.
\]
Then
\begin{equation}\label{eq:est-classical}
 \mu\|u_0\|_{H^2(\Omega)}
 +(\lambda+\mu)\|\operatorname{div}u_0\|_{H^1(\Omega)}
 \le C_\Omega\|f\|_{L^2(\Omega)}.
\end{equation}
\end{theorem}

Set $\beta=(\lambda+\mu)/\mu$, $\varepsilon=\beta^{-1}$,
$\pi=-\beta\operatorname{div}u_0$ and $F=f/\mu$. The normalized system is
\begin{equation}\label{eq:est-classical-mixed}
 -\Delta u+\nabla\pi=F,
 \qquad \operatorname{div}u+\varepsilon\pi=0,
 \qquad u\in H_0^1(\Omega)^3,
 \qquad \pi\in L_0^2(\Omega).
\end{equation}
We include $\varepsilon=0$, which is the incompressible Stokes system.
The usual divergence inf--sup condition and the energy identity give
\begin{equation}\label{eq:est-classical-energy}
 \|u\|_{H^1}+\|\pi\|_{L^2}
 \le C_\Omega\|F\|_{H^{-1}},\qquad 0\le\varepsilon\le1.
\end{equation}
Indeed, testing with $u$ controls $\|\nabla u\|$ and
$\varepsilon^{1/2}\|\pi\|$; the first equation and the inf--sup condition
then control the unweighted pressure norm. For $\varepsilon>0$ its mean
vanishes by integration of the constraint. The argument never divides an
energy bound by $\varepsilon$.

\subsection{A planar estimate with the correct constraint space}

\begin{lemma}\label{lem:est-planar}
Let $D\subset\mathbb R^2$ be a fixed bounded convex polygon. Suppose
$0\le\varepsilon\le1$, $U\in H_0^1(D)^2$ and $P\in L^2(D)$ satisfy
\[
 -\Delta U+\nabla P=F\in L^2(D)^2,
 \qquad \operatorname{div}U+\varepsilon P=g\in H_0^1(D).
\]
No mean condition is imposed on $P$. Then
\begin{equation}\label{eq:est-planar}
 \|U\|_{H^2(D)}+\|P\|_{H^1(D)}
 \le C_D\bigl(\|F\|_{L^2(D)}+\|g\|_{H^1(D)}+\|P\|_{L^2(D)}\bigr),
\end{equation}
uniformly for $0\le\varepsilon\le1$.
\end{lemma}

\begin{proof}
Choose $\phi\in C_c^\infty(D)$ with $\int_D\phi=1$, and put
\[
 m=\int_D P,\qquad \widetilde P=P-m\phi,
 \qquad \widetilde g=g-\varepsilon m\phi.
\]
The integrated constraint gives $\varepsilon m=\int_Dg$.
Consequently $\widetilde g\in H_0^1(D)\cap L_0^2(D)$ and
$\|\widetilde g\|_{H^1}\le C_D\|g\|_{H^1}$.
The support-preserving Bogovskii operator
\cite{CostabelMcIntosh2010} provides
$w\in H_0^2(D)^2$ such that
\[
 \operatorname{div}w=\widetilde g,
 \qquad \|w\|_{H^2}\le C_D\|g\|_{H^1}.
\]
For $v=U-w$ one obtains
\[
 -\Delta v+\nabla\widetilde P
 =F+\Delta w-m\nabla\phi=:H,
 \qquad \operatorname{div}v+\varepsilon\widetilde P=0,
 \qquad \int_D\widetilde P=0.
\]
If $\varepsilon>0$, elimination of $\widetilde P$ gives
$-\Delta v-\varepsilon^{-1}\nabla\operatorname{div}v=H$.
The planar uniform elasticity theorem of
B\u{a}cu\c{t}\u{a}--Bramble \cite[Theorem 3.1]{BacutaBramble2003}
implies $\|v\|_{H^2}\le C_D\|H\|_{L^2}$.
The equation and the zero pressure mean then yield
$\|\widetilde P\|_{H^1}\le C_D\|H\|_{L^2}$.
At $\varepsilon=0$, the same estimate follows by solving the penalized
systems for $\varepsilon_n\downarrow0$, using these uniform bounds, and
identifying their weak limit by Stokes energy uniqueness.
Finally $|m|\le |D|^{1/2}\|P\|_{L^2}$, which proves
\eqref{eq:est-planar}.
\end{proof}

The bump correction is essential to the stated proof: subtracting a
constant pressure mean would generally destroy the zero trace of the
constraint datum. No right inverse on arbitrary $H^1(D)$ divergence data
is used here.

\subsection{Straight edges and all axial frequencies}

Let $W\subset\mathbb R^2$ be a fixed convex wedge. At axial Fourier
frequency $\sigma\in\{-1,1\}$, set
\[
 D_\sigma=(\partial_1,\partial_2,i\sigma),\qquad
 B_\sigma U=\operatorname{div}_\perp U_\perp+i\sigma U_3.
\]
Consider
\begin{equation}\label{eq:est-edge-unit}
 - (\Delta_\perp-1)U+D_\sigma P=F,
 \qquad B_\sigma U+\varepsilon P=G.
\end{equation}
The map $B_\sigma:H_0^1(W)^3\to L^2(W)$ has a bounded right inverse.
To see this, cover $W$ by fixed-size convex polygonal patches with bounded
overlap and finitely many shapes, and use a subordinate partition
$\chi_j$. For $h_j=\chi_jh$, put $c_j=\int h_j$ and choose a unit-integral
interior bump $\phi_j$. Lift $h_j-c_j\phi_j$ by the ordinary planar
Bogovskii operator and lift $c_j\phi_j$ by the axial component
$c_j\phi_j/(i\sigma)$. Zero extension and summation give
\[
 B_\sigma R_\sigma h=h,
 \qquad \|R_\sigma h\|_{H^1(W)}\le C_W\|h\|_{L^2(W)}.
\]
In particular no total integral of $h$ on the unbounded wedge is required.
The velocity form associated with $-\Delta_\perp+1$ is coercive on
$H_0^1(W)^3$. Its pressure Schur complement is coercive on the whole
$L^2(W)$ space. Adding $\varepsilon I$ preserves this property, so
\begin{equation}\label{eq:est-edge-energy}
 \|U\|_{H^1(W)}+\|P\|_{L^2(W)}
 \le C_W\bigl(\|F\|_{H^{-1}(W)}+\|G\|_{L^2(W)}\bigr)
\end{equation}
uniformly for $0\le\varepsilon\le1$.

Write $U=(v,w)$. Its planar equations have right sides
$F_\perp-v$ and $G-i\sigma w$; the latter belongs to $H_0^1(W)$ whenever
$G\in H_0^1(W)$. After a patch cutoff, the additional constraint term is
$\nabla\chi_j\cdot v$, which also has zero trace. Thus
Lemma~\ref{lem:est-planar} applies to the planar components.
The axial equation is the scalar Poisson equation
$-\Delta_\perp w=F_3-w-i\sigma P$.
Convex-polygon Poisson regularity, bounded overlap, and
\eqref{eq:est-edge-energy} give
\begin{equation}\label{eq:est-edge-strong-unit}
 \|U\|_{H^2(W)}+\|P\|_{H^1(W)}
 \le C_W\bigl(\|F\|_{L^2(W)}+\|G\|_{H^1(W)}\bigr).
\end{equation}

For any nonzero axial frequency $\xi$, write $\tau=|\xi|$ and $X=\tau x$.
The substitutions
\[
 \widetilde U(X)=U(X/\tau),\quad
 \widetilde P(X)=\tau^{-1}P(X/\tau),\quad
 \widetilde F(X)=\tau^{-2}F(X/\tau),\quad
 \widetilde G(X)=\tau^{-1}G(X/\tau)
\]
reduce the problem to \eqref{eq:est-edge-unit}, without changing
$\varepsilon$. With
$\|v\|_{k,\tau}^2=\sum_{j=0}^k\tau^{2(k-j)}|v|_{H^j}^2$,
the common scaling factor in \eqref{eq:est-edge-strong-unit} cancels,
giving
\[
 \|U\|_{2,\tau}+\|P\|_{1,\tau}
 \le C_W(\|F\|_{L^2}+\|G\|_{1,\tau}).
\]
Use this bound for $|\xi|\ge1$. For $|\xi|\le1$, the unscaled patch
argument retains the energy remainder
$\|U\|_{H^1}+\|P\|_{L^2}$, and all frequency coefficients are bounded.
Plancherel's theorem therefore proves the three-dimensional edge estimate
\begin{equation}\label{eq:est-classical-edge}
 \|u\|_{H^2(W\times\mathbb R)}+\|p\|_{H^1(W\times\mathbb R)}
 \le C_W\bigl(\|F\|_{L^2}+\|G\|_{H^1}
                   +\|u\|_{H^1}+\|p\|_{L^2}\bigr),
\end{equation}
for $G\in H_0^1$. The half-space and interior estimates follow by the
same argument.

\subsection{A uniform classical vertex pencil}

The spherical section $\Sigma$ of a genuine convex vertex cone has closure
in an open hemisphere $H$. Zero extension maps $H_0^1(\Sigma)$ into
$H_0^1(H)$. We use the hemisphere inequalities
\begin{equation}\label{eq:est-hemisphere}
 \|\nabla_Sv\|^2\ge2\|v\|^2,
 \qquad
 \|\nabla_Sv\|^2\ge6\left(\|v\|^2-
                 |H|^{-1}\left|\int_Hv\right|^2\right).
\end{equation}
The second is the hemisphere variance inequality in
\cite[Lemma 2.3, equation (2.25)]{MayborodaMazya2008}, extended from
$H_0^2$ to $H_0^1$ by density. It can also be obtained by separating the
azimuthal mean, whose Neumann variance gap is $6$, from the remaining
modes, whose first Dirichlet eigenvalue is at least $6$.

Set
\[
 s=-\tfrac12+x+iy,\quad 0\le x\le1,\quad
 t=-1-\overline s,\quad Y=y^2,\quad \kappa=(1+y^2)^{1/2}.
\]
For $U\in H_0^1(\Sigma)^3$, write
$R=U\cdot\omega$, $W=U-R\omega$, and
$d_sU=\operatorname{div}_SW+(s+2)R=:h$.
The normalized Lam\'e pencil is
\[
 b_{\beta,s}(U,V)
 = (\nabla_SU,\nabla_SV)-s(s+1)(U,V)
       +\beta(d_sU,d_tV).
\]
Put $A=\|\nabla_SU\|^2$, $N=\|U\|^2$, and use subscripts $R$ for
the corresponding scalar norms. Integration by parts gives
\[
 b_{0,s}(U,R\omega)
 =A_R-(s+1)(s+2)N_R+2(h,R).
\]
For $\delta=2/\beta$, define
\[
 T_{s,\beta}U=W+
       \frac{s+2+\delta}{1-\overline s+\delta}R\omega.
\]
Both $T_{s,\beta}$ and its inverse are uniformly bounded on the parameter
space $V_y=H_0^1(\Sigma)^3$ with norm
$\|U\|_{V_y}^2=A+\kappa^2N$. Indeed, the radial denominator has real
part at least $1/2$, its modulus ratio is at most $5$, and the radial
projection is a fixed smooth multiplier. A direct calculation yields
\begin{equation}\label{eq:est-classical-cancel}
 b_{\beta,s}(U,T_{s,\beta}U)
 =B_0+\frac{2x}{1-s+\delta}B_R^0+\beta\|h\|^2,
\end{equation}
where
$B_0=A-s(s+1)N$ and
$B_R^0=A_R-(s+1)(s+2)N_R$.

The integral of the tangential divergence vanishes, so
$\int_Hh=(s+2)\int_HR$, after zero extension. Consequently
\eqref{eq:est-hemisphere} implies
\[
 A_R\ge6N_R-\frac{6}{|s+2|^2}\|h\|^2.
\]
Set $d=3/2-x+\delta$ and
$\alpha=(x+1/2)(x+3/2)\le15/4$. The exact real numerator is
\[
 (d^2+Y)\operatorname{Re}\frac{B_R^0}{1-s+\delta}
 =d(A_R-\alpha N_R)+Y(7/2+x+\delta)N_R.
\]
All terms except the mean contribution are nonnegative. Its loss is at
most $12x\|h\|^2/(d|s+2|^2)$, and
\[
 \frac{12x}{\beta d|s+2|^2}
 \le \frac{12x}{[\beta(3/2-x)+2](x+3/2)^2}
 \le\frac{96}{125}.
\]
Here $x/(x+3/2)^2\le4/25$ and
$\beta(3/2-x)+2\ge5/2$. Also
$\operatorname{Re}B_0\ge(5/8)A+YN$. Thus
\begin{equation}\label{eq:est-classical-coercive}
 \operatorname{Re}b_{\beta,s}(U,T_{s,\beta}U)
 \ge\frac58 A+YN+\frac{29}{125}\beta\|d_sU\|^2.
\end{equation}
This holds for every $\beta\ge1$, every $0\le x\le1$, and every
$y\in\mathbb R$. In particular it controls $\|U\|_{V_y}^2$ uniformly.

For completeness, the mixed pencil is
\begin{equation}\label{eq:est-classical-mixed-pencil}
 b_{0,s}(U,V)-(P,d_tV)=\mathcal F(V),
 \qquad d_sU+\varepsilon P=\mathcal G.
\end{equation}
The maps $d_s,d_t:V_y\to L^2(\Sigma)$ have uniform right inverses.
One construction uses a smooth gnomonic chart to a fixed convex polygon,
patches of diameter $\eta\asymp\kappa^{-1}$, and surface Piola transforms.
For $h_j=\chi_jh$, let $c_j=\int h_j$ and choose an interior unit-integral
bump $\phi_j$. A tangential Bogovskii lift of $h_j-c_j\phi_j$ satisfies
\[
 \|W_j\|_{L^2}\le C\eta\|h_j\|_{L^2},
 \qquad \|\nabla_SW_j\|_{L^2}\le C\|h_j\|_{L^2}.
\]
Add $c_j\phi_j\omega/(s+2)$. Since
$|c_j|\le C\eta\|h_j\|$, $\|\phi_j\|\le C\eta^{-1}$,
$\|\nabla_S\phi_j\|\le C\eta^{-2}$, and
$|s+2|\asymp\kappa$, the added field has the same parameter norm bound.
Summing over bounded-overlap patches gives the right inverse of $d_s$.
The construction for $d_t$ uses $t+2=1-\overline s$, whose modulus is
also comparable to $\kappa$ throughout the strip.

At $\varepsilon=0$, let $\beta\to\infty$ in the test multiplier.
The limit $T_s$ maps $\ker d_s$ isomorphically onto $\ker d_t$.
Equation~\eqref{eq:est-classical-coercive}, with $d_sU=0$, solves the
kernel problem. The two divergence right inverses lift $\mathcal G$ and
recover $P$ from the momentum residual. Hence the full mixed inverse obeys
\begin{equation}\label{eq:est-classical-mixed-inverse}
 \|U\|_{V_y}+\|P\|_{L^2}
 \le C_\Sigma(\|\mathcal F\|_{V_y^*}+\|\mathcal G\|_{L^2}).
\end{equation}
The perturbation $(U,P)\mapsto(0,\varepsilon P)$ is bounded on precisely
these weak mixed spaces. A uniform Neumann series handles sufficiently
small $\varepsilon$. For the remaining compact interval, eliminate
$P=\varepsilon^{-1}(\mathcal G-d_sU)$ and apply
\eqref{eq:est-classical-coercive}, with bounded
$\beta=\varepsilon^{-1}$. This proves
\eqref{eq:est-classical-mixed-inverse} for all $0\le\varepsilon\le1$.
The operator family and its inverse are holomorphic in $s$ on the strip.
No perturbation is asserted on a strong constraint space $H_0^1$.

\subsection{Mellin identification and the global estimate}

For a radial vertex cutoff $\chi$, put $v=\chi u$ and $p=\chi\pi$,
and extend them by zero to the infinite vertex cone $K$. The localized
equations have data
\[
 F_\chi=\chi F-2\nabla\chi\cdot\nabla u-(\Delta\chi)u+\pi\nabla\chi,
 \qquad G_\chi=\nabla\chi\cdot u\in H_0^1(K),
\]
and
\begin{equation}\label{eq:est-classical-local-data}
 \|F_\chi\|_{L^2}+\|G_\chi\|_{H^1}
 \le C_\chi(\|F\|_{L^2}+\|u\|_{H^1}+\|\pi\|_{L^2}).
\end{equation}
The zero trace of $G_\chi$ follows from $u=0$ on the physical faces and
from the cutoff at artificial boundaries; it does not require $u\in H^2$.

Write $t=\log r$. At the energy line $\operatorname{Re}s=-1/2$,
\[
 U_0=r^{1/2}v\in H_0^1(\mathbb R\times\Sigma)^3,
 \qquad P_0=r^{3/2}p\in L^2(\mathbb R\times\Sigma).
\]
Hardy's inequality $\|v/r\|_{L^2}\le2\|\nabla v\|_{L^2}$ proves
this assertion directly from energy regularity. The target line
$\operatorname{Re}s=1/2$ has data
\[
 F_1=r^{3/2}F_\chi,
 \qquad G_1=r^{1/2}G_\chi,
 \qquad
 \|F_1\|_{L^2_{\rm cyl}}=\|F_\chi\|_{L^2},\quad
 \|G_1\|_{L^2_{\rm cyl}}=\|G_\chi/r\|_{L^2}
 \le2\|\nabla G_\chi\|_{L^2}.
\]
The energy-line data equal $e^t(F_1,G_1)$ and have a common upper
support bound in $t$. As in the proof of Lemma~\ref{lem:g-mellin-inversion}, approximate
$(F_1,G_1)$ by
$C_c^\infty(\mathbb R;L^2(\Sigma)^3\times L^2(\Sigma))$ data
with a common upper support bound. The approximation converges in the
data spaces on both lines. For each approximant,
holomorphy of the full inverse
\eqref{eq:est-classical-mixed-inverse} and rapid vertical decay of the
transformed data justify contour shifting. The line solutions obey
$(U_{1,n},P_{1,n})=e^{-t}(U_{0,n},P_{0,n})$.
Both sequences converge by the line-inverse bounds. Passing through the
multiplier $e^{-t}$ only on bounded $t$ intervals identifies the target
solution, as a distribution, with
$(r^{-1/2}v,r^{1/2}p)$. The initial solution is the actual energy solution
by uniqueness on the energy line. Consequently
\begin{equation}\label{eq:est-classical-weights}
 \|Dv/r\|_{L^2}+\|v/r^2\|_{L^2}+\|p/r\|_{L^2}
 \le C\bigl(\|F_\chi\|_{L^2}+\|G_\chi/r\|_{L^2}\bigr).
\end{equation}
This argument constructs the target-space solution before identifying it;
it never assumes the desired vertex weights for the energy solution.

On fixed nested vertex annuli $A\Subset A^*$, a finite interior/face/edge
cover and \eqref{eq:est-classical-edge} give a strong local estimate.
After rescaling an annulus of radius $h$, it reads
\begin{align*}
 \|D^2v\|_{L^2(hA)}+\|Dp\|_{L^2(hA)}
 \le C\bigl(&\|F_\chi\|_{L^2(hA^*)}+\|DG_\chi\|_{L^2(hA^*)}
       +h^{-1}\|G_\chi\|_{L^2(hA^*)}\\
 &+h^{-1}\|Dv\|_{L^2(hA^*)}+h^{-2}\|v\|_{L^2(hA^*)}
       +h^{-1}\|p\|_{L^2(hA^*)}\bigr).
\end{align*}
Square and sum over dyadic annuli. Finite overlap,
\eqref{eq:est-classical-weights}, and Hardy's inequality for $G_\chi$
control all terms. A finite cover of the remaining faces, edges and
interior gives
\[
 \|u\|_{H^2(\Omega)}+\|\pi\|_{H^1(\Omega)}
 \le C_\Omega(\|F\|_{L^2}+\|u\|_{H^1}+\|\pi\|_{L^2}).
\]
Use \eqref{eq:est-classical-energy} and restore $\mu$ to obtain
Theorem~\ref{thm:est-classical}.

\section{Strain-gradient estimates as a regularity consequence}
\label{sec:est-sge}

The three-dimensional regularity results now supply all inputs of the
abstract SGE estimate in~\cite[Theorem~3.3]{PolygonCompanion2026}.
We record the resulting bounds and the short parameter-scaling argument.

\begin{proposition}\label{prop:est-sge-transfer}
Let $\Omega\subset\mathbb R^3$ be a fixed bounded convex polyhedron.
Let $u_\iota\in H_0^2(\Omega)^3$ and $u_0\in H_0^1(\Omega)^3$ solve
\[
 (I-\iota^2\Delta)\mathcal L_{\lambda,\mu}u_\iota=f,
 \qquad \mathcal L_{\lambda,\mu}u_0=f,
 \qquad
 \mathcal L_{\lambda,\mu}u=-\mu\Delta u
                     -(\lambda+\mu)\nabla\operatorname{div}u.
\]
Set $p_\iota=\lambda\operatorname{div}u_\iota$ and
$p_0=\lambda\operatorname{div}u_0$. Then, for all
$\lambda\ge0$, $\mu>0$ and $0<\iota<1$,
\begin{align}
 |u_\iota-u_0|_{H^1}+\iota\|u_\iota\|_{H^2}
                  +\iota^2\|u_\iota\|_{H^3}
 &\le C_\Omega\mu^{-1}\sqrt\iota\,\|f\|_{L^2},
 \label{eq:est-sge-u}\\
 \|p_\iota-p_0\|_{L^2}+\iota|p_\iota|_{H^1}
                  +\iota^2\|p_\iota\|_{H^2}
 &\le C_\Omega\frac{\lambda}{\lambda+2\mu}
                    \sqrt\iota\,\|f\|_{L^2}.
 \label{eq:est-sge-p}
\end{align}
\end{proposition}

\begin{proof}
For \(\mu=1\), Theorems~\ref{thm:g-stokes},
\ref{thm:g-lameuniform} and~\ref{thm:est-classical} verify, respectively,
Assumptions~2.4, 2.5 and estimate~(3.13) of
\cite{PolygonCompanion2026}. Theorem~3.3 of that paper is formulated
in arbitrary dimension and therefore gives the six-term bound here.
The same argument yields the residual estimate
\begin{equation}\label{eq:est-sge-third}
 \iota^2\|u_\iota\|_{H^3}\le C|u_\iota-u_0|_{H^1},\qquad \mu=1,
\end{equation}
by its Lemma~3.2; this will also be used in the boundary-trace criterion.

For completeness, the sharper pressure factor follows without division
by \(\lambda\). Set
\[
 D=\|\operatorname{div}(u_\iota-u_0)\|_{L^2}
 +\iota|\operatorname{div}u_\iota|_{H^1}
 +\iota^2\|\operatorname{div}u_\iota\|_{H^2}.
\]
The displacement and pressure estimates imply
\(2D\le C\sqrt\iota\|f\|_{L^2}\) and
\(\lambda D\le C\sqrt\iota\|f\|_{L^2}\).
Thus \((\lambda+2)D\le C\sqrt\iota\|f\|_{L^2}\), and multiplying
by \(\lambda/(\lambda+2)\) proves the refined pressure bound.
Finally, for the same load and \(\kappa=\lambda/\mu\),
\[
 u_{\iota;\lambda,\mu}=\mu^{-1}u_{\iota;\kappa,1},\qquad
 p_{\iota;\lambda,\mu}=p_{\iota;\kappa,1},
\]
with identical scaling for the classical solution. This proves the
stated bounds for every \(\mu>0\). The material factors and this
abstract transfer are shared with the companion paper; the new inputs
are the polyhedral regularity theorems established above.
\end{proof}

\section{Simultaneous sharpness for one fixed load}
\label{sec:est-sharpness}

\begin{theorem}\label{thm:est-sharpness}
Let $\Omega\subset\mathbb R^3$ be a fixed bounded convex polyhedron.
There is a nonzero polynomial load $f_\Omega$, independent of
$\lambda,\mu,\iota$, for which all six individual quantities in
\eqref{eq:est-sge-u}--\eqref{eq:est-sge-p} attain their comparison scales. More precisely, with
\[
\begin{array}{lll}
 T_1=|u_\iota-u_0|_{H^1},&T_2=\iota\|u_\iota\|_{H^2},
       &T_3=\iota^2\|u_\iota\|_{H^3},\\
 T_4=\|p_\iota-p_0\|_{L^2},&T_5=\iota|p_\iota|_{H^1},
       &T_6=\iota^2\|p_\iota\|_{H^2},
\end{array}
\]
there are $c_\Omega,C_\Omega>0$ such that, for every
$\lambda\ge0$, $\mu>0$ and $0<\iota<1$,
\[
 c_\Omega\mu^{-1}\sqrt\iota\le T_j
       \le C_\Omega\mu^{-1}\sqrt\iota\quad(j=1,2,3),
\]
\[
 c_\Omega\alpha\sqrt\iota\le T_j
       \le C_\Omega\alpha\sqrt\iota\quad(j=4,5,6),
 \qquad \alpha=\frac{\lambda}{\lambda+2\mu}.
\]
At $\lambda=0$, all three pressure quantities and their comparison
scales are zero.
\end{theorem}

\begin{proof}
Use the genuine maximal planar faces of $\Omega$, merging coplanar
subdivisions. For each face $F$, let $\ell_F$ be its unit affine distance,
positive in $\Omega$. Fix $a=e_3$ and define
\begin{equation}\label{eq:est-fixed-load}
 \psi=\prod_F\ell_F^2,
 \qquad w=\nabla\psi\times a,
 \qquad z=\nabla\psi,
 \qquad f_\Omega=-\Delta w-\nabla\Delta\psi.
\end{equation}
These are polynomials depending only on the fixed domain. Since
$\psi=\nabla\psi=0$ on $\partial\Omega$, one has
$\psi\in H_0^2$, $w,z\in H_0^1{}^3$, and
$\operatorname{div}w=0$. Direct substitution shows that the exact
classical solution is
\begin{equation}\label{eq:est-fixed-classical}
 u_0=\mu^{-1}w+(\lambda+2\mu)^{-1}z,
 \qquad p_0=\alpha\Delta\psi.
\end{equation}

In the relative interior of a face $F$, with outward normal $n_F$,
\[
 D^2\psi=h_F n_F\otimes n_F,
 \qquad h_F=2\prod_{E\ne F}\ell_E^2>0.
\]
Writing $h=h_F$ facewise gives
\[
 \operatorname{Tr}\Delta\psi=h,
 \qquad \partial_nz=hn,
 \qquad \partial_nw=h(n\times a).
\]
Thus
\[
 B_p=\|h\|_{L^2(\partial\Omega)}>0,
 \qquad
 B_w^2=\sum_F|n_F\times a|^2\|h_F\|_{L^2(F)}^2>0.
\]
For the second inequality, boundedness of the polyhedron implies that
its face normals span $\mathbb R^3$, so they cannot all be parallel to
$a$. The two terms of
\[
 \partial_nu_0=h\bigl(\mu^{-1}(n\times a)
                      +(\lambda+2\mu)^{-1}n\bigr)
\]
are orthogonal on each face. Also
$(f_\Omega,w)=\|\nabla w\|_{L^2}^2>0$, so the fixed load is nonzero.

We first set $\mu=1$. Then $\|u_0\|_{H^2}\le M_u$ and
$\|\partial_nu_0\|_{L^2(\partial\Omega)}\ge B_w$, with $M_u,B_w$
independent of $\lambda$. Introduce the pressure surrogate
\[
 q_\iota=(\lambda+2)\operatorname{div}u_\iota,
 \qquad q_0=\Delta\psi,
\]
and set
\[
 x=|u_\iota-u_0|_{H^1},\quad y=\iota\|u_\iota\|_{H^2},
 \quad z_*=\iota^2\|u_\iota\|_{H^3},
\]
\[
 X=\|q_\iota-q_0\|_{L^2},\quad
 Y=\iota|q_\iota|_{H^1},\quad
 Z=\iota^2\|q_\iota\|_{H^2}.
\]
The upper bounds give
\begin{equation}\label{eq:est-sharp-upper}
 x+y+z_*+X+Y+Z\le C_\Omega\sqrt\iota.
\end{equation}
For the surrogate quantities, this follows by combining unweighted
divergence continuity with the $\lambda$-weighted estimates. It involves
no division by $\lambda$. The clamped condition implies
$q_\iota\in H_0^1$, whereas $\operatorname{Tr}q_0=h$ and
$\|q_0\|_{H^1}=M_q$ are fixed.

The multiplicative trace inequality applied to
$\nabla(u_\iota-u_0)$ and to $q_\iota-q_0$ yields
\begin{equation}\label{eq:est-trace-lower}
 B_w^2\le Cx(\iota^{-1}y+M_u),
 \qquad B_p^2\le CX(\iota^{-1}Y+M_q).
\end{equation}
Inserting the upper bounds for $y,Y$ gives
$x,X\ge c\sqrt\iota$ throughout $0<\iota<1$.
Inserting the upper bounds for $x,X$ instead yields
\[
 y,Y\ge c\sqrt\iota-C\iota.
\]
Thus $y,Y\ge c'\sqrt\iota$ below a fixed threshold independent of
$\lambda$.

We use the shifted interpolation inequality
\begin{equation}\label{eq:est-shifted-interpolation}
 \|v\|_{H^{m+1}}^2
 \le C\bigl(\|v-w_0\|_{H^m}\|v\|_{H^{m+2}}
                      +\|w_0\|_{H^{m+1}}^2\bigr),\qquad m=0,1.
\end{equation}
This dimension-independent inequality is proved in
\cite[proof of Proposition~3.5]{PolygonCompanion2026} using a common
Sobolev extension operator and a Fourier inner product.
With $(m,v,w_0)=(1,u_\iota,u_0)$ and then
$(0,q_\iota,q_0)$, the lower bounds for $y,Y$ imply
\[
 c\iota^{-1}\le C\sqrt\iota\,\|u_\iota\|_{H^3}+CM_u^2,
 \qquad
 c\iota^{-1}\le C\sqrt\iota\,\|q_\iota\|_{H^2}+CM_q^2.
\]
Hence $z_*,Z\ge c\sqrt\iota$ for
$0<\iota<\iota_*$, with one common parameter-independent
$\iota_*>0$.

To cover the rest of the interval, choose
$\phi_n\in C_c^\infty(\Omega)$ with $\phi_n\to\psi$ in $H^2$, and
put $v_n=\nabla\phi_n\times a$. Then
\[
 (f_\Omega,v_n)\longrightarrow\|\nabla w\|_{L^2}^2>0,
 \qquad
 (\Delta^2\psi,\phi_n)\longrightarrow\|\Delta\psi\|_{L^2}^2>0.
\]
Fix one sufficiently large $n$, and write $\phi=\phi_n$, $v=v_n$.
The vector $v$ is compactly supported and divergence free. Testing the
SGE equation against it eliminates every $\lambda$ term and gives
\[
 0<|(f_\Omega,v)|
 =|(u_\iota,\iota^2\Delta^2v-\Delta v)|
 \le C_v\|u_\iota\|_{L^2},\qquad 0<\iota<1.
\]
Moreover $\operatorname{div}f_\Omega=-\Delta^2\psi$, so
\[
 (I-\iota^2\Delta)\Delta q_\iota=\Delta^2\psi
\]
in distributions. Testing with the fixed $\phi$ gives
\[
 0<|(\Delta^2\psi,\phi)|
 \le(\|\Delta\phi\|_{L^2}+\|\Delta^2\phi\|_{L^2})
                  \|q_\iota\|_{L^2}.
\]
Since $q_\iota\in H_0^1$, Poincar\'e's inequality now yields
\[
 \|u_\iota\|_{H^2},\ \|u_\iota\|_{H^3},\
 |q_\iota|_{H^1},\ \|q_\iota\|_{H^2}\ge c_\Omega>0
\]
uniformly for every $\lambda\ge0$ and $0<\iota<1$.
For $\iota_*\le\iota<1$ it follows that
\[
 y,Y\ge c\iota\ge c\sqrt{\iota_*}\sqrt\iota,
 \qquad
 z_*,Z\ge c\iota^2\ge c\iota_*^{3/2}\sqrt\iota.
\]
Together with \eqref{eq:est-trace-lower}, this proves all six surrogate
and displacement lower bounds on the full interval.

Finally $p_\iota=\alpha q_\iota$ and $p_0=\alpha q_0$ when $\mu=1$.
Multiply the three surrogate bounds by $\alpha=\lambda/(\lambda+2)$.
This includes $\lambda=0$ exactly. The material scaling used in
Proposition~\ref{prop:est-sge-transfer}, applied to the same fixed
$f_\Omega$, gives the factor $\mu^{-1}$ in displacement and replaces
$\alpha$ by $\lambda/(\lambda+2\mu)$ in pressure.
\end{proof}

On the unit cube, one may take
$\psi=\prod_{j=1}^3x_j^2(1-x_j)^2$.
With $a=e_3$, exact polynomial integration gives
$B_p^2=6/99225$ and $B_w^2=4/99225$.
The general proof only uses their strict positivity and applies to every
fixed bounded convex polyhedron, without an additional vertex-angle or
spherical-eigenvalue threshold.


\subsection{When the half-order displacement rate is active}

We write $a_\iota=\Theta(b_\iota)$ when $c\,b_\iota\le a_\iota
\le C\,b_\iota$ for all sufficiently small $\iota>0$, with constants
$c,C>0$ at the stated fixed parameters and load.

\begin{proposition}[A boundary-trace criterion]\label{prop:boundary-active}
Fix a bounded convex polyhedron, material parameters \(\lambda\ge0,\mu>0\),
and a load \(f\in L^2(\Omega)^3\). Each of \(T_1,T_2,T_3\) is
\(\Theta(\sqrt{\iota})\) as \(\iota\downarrow0\) if and only if
\[
 \|\partial_nu_0\|_{L^2(\partial\Omega)}>0.
\]
If this trace vanishes, all three quantities are \(O(\iota)\).
The comparison constants in this proposition may depend on the fixed
parameters and load.
\end{proposition}
\begin{proof}
The sufficiency argument in
\cite[Proposition~3.5]{PolygonCompanion2026} uses only the
multiplicative trace inequality, the shifted interpolation inequality
and the three displacement upper bounds. These are valid here by
Theorem~\ref{thm:est-classical}, Proposition~\ref{prop:est-sge-transfer}
and~\eqref{eq:est-shifted-interpolation}.
If the normal trace vanishes, \(u_0\in H^2_0\), and energy comparison
gives \(T_1+T_2=O(\iota)\). Estimate~\eqref{eq:est-sge-third} and
material scaling give \(T_3=O(\iota)\). This is the same
boundary-trace argument as in the companion paper, with its regularity
inputs now established on convex polyhedra.
\end{proof}



\section{Necessity of convexity}\label{sec:necessity}
Throughout this section, $\Omega$ is a bounded connected polyhedral
Lipschitz domain in $\mathbb R^3$.  In particular, its boundary consists of
finitely many planar faces; coplanar subdivisions are not counted as
geometric edges.

\begin{lemma}[A nonconvex polyhedron has a reentrant edge]\label{lem:reentrant}
If $\Omega$ is not convex, it has an edge with interior dihedral angle
$\omega\in(\pi,2\pi)$.
\end{lemma}
\begin{proof}
Suppose that every interior dihedral angle is at most $\pi$.
The closed set $K=\overline\Omega$ is locally convex at points in the
interiors of faces and edges.  At a vertex, the Lipschitz property gives,
after a rigid coordinate change, a local representation of $\Omega$ as
the epigraph of a continuous function $f$ on a small convex planar disk.
Since the boundary has finitely many planar faces, $f$ is piecewise affine.
The dihedral-angle assumption says that each crease is convex.  On a
generic line segment avoiding partition vertices and overlapping creases,
the slope of the restriction of $f$ is nondecreasing across each crease.
Its restriction is therefore convex.  Approximation of arbitrary segments
by such segments, together with continuity, proves that $f$ is convex.
Thus $K$ is locally convex also at vertices.

The Tietze--Nakajima theorem states that a closed connected locally convex
subset of Euclidean space is convex; see \cite[Theorem 1]{BjorndahlKarshon2010}.
It applies to $K$.  A Lipschitz domain is the interior of its closure, so
$\Omega$ is convex, a contradiction.
\end{proof}

Choose an interior point of a reentrant edge and use $z$ as the coordinate
along that edge.  In a sufficiently small cylinder, $\Omega$ agrees with
\[
 \{(r,\theta,z):0<r<r_*,\quad 0<\theta<\omega,\quad |z|<z_*\}.
\]
The functions below vanish before reaching the artificial radial and axial
boundaries.  Consequently, their zero extensions define functions on
$\Omega$; no global prism structure is required.

\begin{lemma}[A clamped scalar edge mode]\label{lem:scalar-edge}
For every $\omega\in(\pi,2\pi)$ there is a nonzero real function
$w(r,\theta)=r^s\Phi(\theta)$ such that
\[
 \Delta_\perp^2w=0,\qquad
 \Phi(0)=\Phi'(0)=\Phi(\omega)=\Phi'(\omega)=0,\qquad
 \frac32<s<1+\frac{\pi}{\omega}<2.
\]
After a radial cutoff, $w$ belongs to $H^2_0\cap W^{2,4}$ but not to $H^3$.
\end{lemma}
\begin{proof}
Put $f_\omega(t)=\sin(t\omega)-t|\sin\omega|$.  Then
\[
 f_\omega(1/2)
 =\sin(\omega/2)\bigl(1+\cos(\omega/2)\bigr)>0,\qquad
 f_\omega(\pi/\omega)=-(\pi/\omega)|\sin\omega|<0.
\]
There is therefore a root $t\in(1/2,\pi/\omega)$; strict concavity on
$(0,\pi/\omega)$ also gives its uniqueness among positive roots in that
interval.  Set $s=1+t$ and
\[
 A_t(\theta)=\sin(t\theta)\sin\theta,\qquad
 B_t(\theta)=t\cos(t\theta)\sin\theta-\sin(t\theta)\cos\theta .
\]
Both functions and their first derivatives vanish at zero.  They are
combinations of the frequencies $s$ and $s-2$, and
\[
 A_t(\omega)B_t'(\omega)-A_t'(\omega)B_t(\omega)
 =\sin^2(t\omega)-t^2\sin^2\omega=0.
\]
Thus
\[
 \Phi(\theta)=B_t(\omega)A_t(\theta)-A_t(\omega)B_t(\theta)
\]
satisfies the four clamped conditions and
$[\partial_\theta^2+(s-2)^2][\partial_\theta^2+s^2]\Phi=0$.
It is nonzero: $A_t(\omega)=-t\sin^2\omega\ne0$, and $A_t,B_t$
are linearly independent, as their Taylor expansions at zero show.
This also verifies directly that the characteristic root is not spurious.
The scalar characteristic equation is classical; see
\cite[(4.3)--(4.5)]{BacutaBramblePasciak2002}.

The estimates $D^kw=O(r^{s-k})$ imply $D^2w\in L^4$ because $s>3/2$.
On the other hand,
\[
 \partial_r^3w=s(s-1)(s-2)r^{s-3}\Phi,\qquad
 \int_0^\epsilon r^{2s-5}\,dr=\infty ,
\]
so $w\notin H^3$.  The clamped traces and an additional cutoff at $r=\epsilon$
give approximation in $H^2$ with error $O(\epsilon^{s-1})$.
The cutoff function is therefore in $H^2_0$.  A finite axial tensor product
does not change these transverse integrability thresholds.
\end{proof}

\subsection{A correction preserving the integrability of the load}
Write $x=(x_1,x_2)$ for the transverse Cartesian coordinates.  Set
\[
 d=s+2,\qquad c_s=\frac4{s+2},\qquad R_s=\frac{xw}{s+2}.
\]
Euler's identity and the product rule give
\begin{equation}\label{eq:edge-euler}
 \operatorname{div}_\perp R_s=w,\qquad
 \Delta_\perp R_s=\frac{x\Delta_\perp w+2\nabla_\perp w}{s+2},\qquad
 \Delta_\perp^2R_s=c_s\nabla_\perp\Delta_\perp w.
\end{equation}
Choose a smooth radial cutoff $\chi$ equal to one near zero and supported
in $r<r_*$, and define
\[
 R=\chi R_s,\qquad
 h=\operatorname{div}_\perp R
 =\left(\chi+\frac{r\chi'}{s+2}\right)w.
\]
On the truncated sector $D$ these functions satisfy
\begin{equation}\label{eq:edge-spaces}
 R\in H^3(D)^2\cap H^2_0(D)^2,\qquad
 h\in H^2_0(D)\cap W^{2,4}(D),\qquad h\notin H^3(D).
\end{equation}
Indeed, $R_s$ is homogeneous of degree $s+1$; its third derivatives are
square integrable.  All zero and first-order traces vanish on the physical
rays, and the functions vanish near the artificial boundary.
The inner-cutoff errors in $H^2$ are $O(\epsilon^s)$ for $R$ and
$O(\epsilon^{s-1})$ for $h$.  Notice that $H^3_0$ membership of $R$ is
neither asserted nor needed.  Also $\int_Dh=0$.

The complete radial cutoff errors are
\[
 q=\Delta_\perp^2h,\qquad
 E=\Delta_\perp^2R-c_s\nabla_\perp\Delta_\perp h.
\]
By \eqref{eq:edge-euler}, $q$ and $E$ vanish near the edge and are bounded
smooth functions supported in a fixed annulus.  In what follows the most
singular remaining integrable factor is
$\Delta_\perp h=O(r^{s-2})\in L^4(D)$.

\begin{proposition}[Fourth-order Lam\'e and SGE counterexamples]
\label{prop:edge-lame}
For each $\lambda\ge0$, $\mu>0$, and $\iota>0$, a nonconvex $\Omega$
admits loads $F,f\in L^4(\Omega)^3$ whose respective clamped energy solutions
of $P_{\lambda,\mu}u=F$ and $(\mathcal L_{\lambda,\mu}+\iota^2P_{\lambda,\mu})u=f$
do not belong to $H^3(\Omega)^3$.
\end{proposition}
\begin{proof}
Put
\[
 \beta=\frac{\lambda+\mu}{\mu},\qquad
 a=\frac{\beta}{\beta+c_s},\qquad
 b=\beta(1-a)=ac_s .
\]
Choose $\eta\in C_c^\infty((-z_*,z_*))$ equal to one on a nonempty
subinterval, and set
\begin{equation}\label{eq:edge-Ubeta}
 U_\beta=(-a\eta'R,\eta h).
\end{equation}
Then
$\operatorname{div}U_\beta=(1-a)\eta'h$ and
$U_\beta\in H^2_0(\Omega)^3\cap W^{2,4}(\Omega)^3$.
Where $\eta=1$ and $\chi=1$, its third component is $w$.
Consequently $U_\beta\notin H^3$.

Let $\mathcal B_\beta=\Delta^2+\beta\nabla\Delta\operatorname{div}$.
The identity $-ac_s+\beta(1-a)=0$ cancels the nonintegrable
term $\eta'\nabla_\perp\Delta_\perp h$.  Keeping all radial cutoff
errors gives the exact formulas
\begin{equation}\label{eq:edge-force4}
\begin{aligned}
 (\mathcal B_\beta U_\beta)_\perp
 &=-a\eta'E-2a\eta'''\Delta_\perp R-a\eta'''''R
       +b\eta'''\nabla_\perp h,\\
 (\mathcal B_\beta U_\beta)_3
 &=\eta q+(2+b)\eta''\Delta_\perp h+(1+b)\eta''''h .
\end{aligned}
\end{equation}
Every displayed term is in $L^4$.
Without the transverse correction in \eqref{eq:edge-Ubeta}, axial
localization would leave $\beta\eta'\nabla_\perp\Delta_\perp w$,
which need not even belong to $L^2$.

For $\mathcal A_\beta=-\Delta-\beta\nabla\operatorname{div}$ one similarly has
\begin{equation}\label{eq:edge-force2}
\begin{aligned}
 (\mathcal A_\beta U_\beta)_\perp
 &=a\eta'\Delta_\perp R+a\eta'''R-b\eta'\nabla_\perp h,\\
 (\mathcal A_\beta U_\beta)_3
 &=-\eta\Delta_\perp h-(1+b)\eta''h .
\end{aligned}
\end{equation}
These terms also belong to $L^4$.
Thus the loads
\[
 F=\mu\mathcal B_\beta U_\beta,\qquad
 f=\mu(\mathcal A_\beta+\iota^2\mathcal B_\beta)U_\beta
\]
have the asserted integrability.  The identities hold distributionally in
$\Omega$.  Density of $C_c^\infty(\Omega)^3$ in $H^2_0(\Omega)^3$
extends them to the energy formulation.  Coercivity and uniqueness identify
$U_\beta$ with the corresponding energy solutions.
\end{proof}

\begin{proposition}[A zero-divergence fourth-order Stokes counterexample]
\label{prop:edge-stokes}
If $\Omega$ is nonconvex, there are
\[
 F\in L^4(\Omega)^3,\quad
 U\in H^2_0(\Omega)^3\setminus H^3(\Omega)^3,\quad
 p\in H^2(\Omega)\cap H^1_0(\Omega)\cap L^2_0(\Omega)
\]
such that $\Delta^2U+\nabla\Delta p=F$ and $\operatorname{div}U=0$.
\end{proposition}
\begin{proof}
With the same $R,h,\eta$, take
\[
 U=(-\eta'R,\eta h),\qquad p=c_s\eta'h .
\]
The divergence is zero and $U\notin H^3$.
The pressure has the stated Sobolev regularity, and its mean is zero
since $\int\eta'=0$.
The force $\Delta^2U+\nabla\Delta p$ is given by
\eqref{eq:edge-force4} with $a=1$, $b=c_s$, hence belongs to $L^4$.
Energy uniqueness is the one established for the mixed problem.
\end{proof}

\begin{remark}[Quantifiers and the second-order problem]\label{rem:quantifiers}
The loads in Proposition~\ref{prop:edge-lame} may depend on the fixed
material parameters and on $\iota$.
They disprove regularity for all loads at each fixed parameter choice.
They do not assert singularity for every load, or the existence of a
single $L^4$ load producing these counterexamples for all parameters.

The analogous necessity for classical $L^2$-to-$H^2$ elasticity follows
from an even simpler mode.  Put $\nu_e=\pi/\omega\in(1/2,1)$,
$k=\chi r^{\nu_e}\sin(\nu_e\theta)$, and $v=\eta k e_3$.
Then $v\in H^1_0\setminus H^2$, whereas
\[
 (\mathcal L_{\lambda,\mu}v)_\perp
 =-(\lambda+\mu)\eta'\nabla_\perp k,\qquad
 (\mathcal L_{\lambda,\mu}v)_3
 =-\mu\eta\Delta_\perp k-(\lambda+2\mu)\eta''k
\]
belong to $L^2$.
\end{remark}

\begin{proof}[Proof of Theorem~\ref{thm:main-convexity}]
On convex polyhedra, the fourth-order Lam\'e and Stokes regularity theorems
give the asserted shifts.  If $u_\iota$ is the SGE energy solution for
$f\in L^2$, then
\[
 P_{\lambda,\mu}u_\iota
 =\iota^{-2}(f-\mathcal L_{\lambda,\mu}u_\iota)\in L^2,
\]
because $u_\iota\in H^2_0$.  The fourth-order shift gives $u_\iota\in H^3$.
The classical assertion follows from Theorem~\ref{thm:est-classical}.
Conversely, Lemma~\ref{lem:reentrant}, Remark~\ref{rem:quantifiers},
and Propositions
\ref{prop:edge-lame}--\ref{prop:edge-stokes} disprove each asserted
regularity property on a nonconvex polyhedron, even with the stronger
loads stated there.  For fixed coefficients, an everywhere-defined
regularity map is bounded by the closed graph theorem, since convergence
in its target identifies the same energy solution by uniqueness.
Parameter-uniform bounds, however, are supplied by the quantitative
convex-domain theorems, not by this closed graph argument.
\end{proof}



\section{A compatible-data consequence of the companion paper}
\label{comp:section}

The following criterion and strong limits are dimension-independent
results of~\cite[Corollary~3.7 and Proposition~3.9]{PolygonCompanion2026}.
They use energy solvability rather than the polyhedral regularity
proved in this paper, and are recalled here to distinguish compatible
loads from the worst-case boundary-layer rates.

\begin{corollary}[Compatible loads, cited criterion and material scaling]
\label{comp:criterion}
Let \(\Omega\subset\mathbb R^3\) be bounded, connected and Lipschitz.
For fixed \(\lambda\ge0\), \(\mu>0\), and \(f\in L^2(\Omega)^3\),
\[
 |u_\iota-u_0|_{H^1}=O(\iota^2)
 \quad\Longleftrightarrow\quad
 u_0\in H^2_0(\Omega)^3,\qquad \Delta f\in H^{-1}(\Omega)^3.
\]
When these conditions hold, let
\(U\in H^1_0(\Omega)^3\) solve
\(\mathcal L_{\lambda,\mu}U=\Delta f\). Then
\[
 \iota^{-2}(u_\iota-u_0)\to U\ \text{in }H^1,\qquad
 \iota^{-2}(p_\iota-p_0)\to\lambda\operatorname{div}U\ \text{in }L^2,
 \qquad \iota|\iota^{-2}(u_\iota-u_0)|_{H^2}\to0.
\]
For \(u_0\ne0\), \(U\ne0\); the pressure limit may vanish.
Writing \(e=u_\iota-u_0\) and \(r=p_\iota-p_0\), the following bounds
hold uniformly for all \(\lambda\ge0\), \(\mu>0\), and \(0<\iota<1\)
whose data satisfy the compatibility conditions:
\begin{align}
 \mu\bigl(|e|_{H^1}+\iota|e|_{H^2}\bigr)
 &\le C_\Omega\iota^2\|\Delta f\|_{H^{-1}},\\
 \|r\|_{L^2}+\iota|r|_{H^1}
 &\le C_\Omega\frac{\lambda}{\lambda+2\mu}
          \iota^2\|\Delta f\|_{H^{-1}}.
\end{align}
\end{corollary}
\begin{proof}
The criterion and limits are exactly the cited proposition, specialized
to \(d=3\). Apply the cited corollary at \(\mu=1\), then use the
material scaling of Proposition~\ref{prop:est-sge-transfer}.
For the sharper pressure factor, put
\(D=\|\operatorname{div}e\|_{L^2}+\iota|\operatorname{div}e|_{H^1}\).
The displacement and pressure bounds give
\(2\mu D+\lambda D\le C_\Omega\iota^2\|\Delta f\|_{H^{-1}}\).
Multiplication by \(\lambda/(\lambda+2\mu)\) proves the claim,
including \(\lambda=0\). No convexity assumption is used.
\end{proof}

\bibliographystyle{abbrv}
\bibliography{../refs}
\end{document}
